\subsection{克赖因-鲁特曼定理}

引理 3. --- 设 \((a_n)_{n \geqslant 0}\) 是一个正实数序列，使得幂级数

\[
f (z) = \sum_ {n \geqslant 0} a _ {n} z ^ {n}
\]

具有有限的收敛半径 r \textgreater{} 0。设 D ⊂ C 是以 0 为心、以 r 为半径的开圆盘。不存在定义在 C 中 r 的某个开邻域 U 上且在 U ∩ D 上与 f 一致的全纯函数 \(\widetilde{f}\)。

假设这样一个全纯函数 \(\tilde{f}\)（定义在 \(r\) 的一个开邻域 U 上）存在。存在实数 \(s < r\) 和 \(\delta > 0\)，使得 \(s + \delta > r\)，且使得以 \(s\) 为心、以 \(\delta\) 为半径的开圆盘 D' 包含于 U 中。\(\tilde{f}\) 在点 \(s\) 处的幂级数展开（VAR, R1, p. 26, 3.2.1）于是对每个 \(z\)（只要它落在以 0 为心、以 \(\delta\) 为半径的圆盘中）都收敛（VAR, R1, p. 29, 3.3.4）。此级数为

\[
\widetilde {f} _ {s} (z) = \sum_ {n = 0} ^ {+ \infty} \frac {\widetilde {f} ^ {(n)} (s)}{n !} z ^ {n} = \sum_ {n = 0} ^ {+ \infty} \frac {f ^ {(n)} (s)}{n !} z ^ {n}
\]

（VAR, R1, p. 27, 3.2.4）。由于 \(s \in D\) ，我们有

\[
f ^ {(n)} (s) = \sum_ {k = n} ^ {+ \infty} k (k - 1) \dots (k - n + 1) a _ {k} s ^ {k - n}
\]

（VAR, R1, p. 28, 3.2.11）。取 \(z\) 使得 \(0 < z < \delta\) 。由于 \(a_{k} \geqslant 0\) ，我们有

\[
\begin{array}{r l} & {\widetilde {f} _ {s} (z) = \sum_ {n = 0} ^ {+ \infty} \biggl (\sum_ {k = n} ^ {+ \infty} \frac {k (k - 1) \cdots (k - n + 1)}{n !} a _ {k} s ^ {k - n} \biggr) z ^ {n}} \\ & {\qquad = \sum_ {k = 0} ^ {+ \infty} \biggl (a _ {k} \sum_ {n = 0} ^ {k} \frac {k !}{n ! (k - n) !} s ^ {k - n} z ^ {n} \biggr) = \sum_ {k = 0} ^ {+ \infty} a _ {k} (s + z) ^ {k}} \end{array}
\]

（TG, III, p. 40）。当 \(s + z > r\) 时，这个表达式的收敛性与幂级数 \(f\) 的收敛半径等于 \(r\) 这一假设相矛盾。

设 E 是一个实巴拿赫空间。设 C 是 E 中的一个尖凸锥。由 C 生成的向量空间等于 C -- C（EVT, II, p. 12, 推论 1）。特别地，锥 C 在 E 中是全的，当且仅当 C -- C 在 E 中稠密。回顾：当 C ∩ (−C) 仅由 0 组成时，称 C 是突出的（EVT, II, p. 11）。

C 的极集 C° 是由连续线性形式 \(\ell \in E'\) 组成的集合，其中 \(\ell(x) \geqslant 0\) 对所有 \(x \in C\) 成立（EVT, II, p. 47, 命题 4）。由双极定理（EVT, II, p. 48, 定理 1），若 C 是闭尖凸锥，则 C°° = C。

引理 4. --- 设 E 是一个实巴拿赫空间，\(u \in \mathcal{L}(\mathrm{E}_{(\mathbf{C})})\) 是 \(\mathrm{E}_{(\mathbf{C})}\) 的一个非零自同态。设 C 是 E 中的一个全凸锥。则存在 \(x \in \mathrm{C}\) 使得 \(u(x) \neq 0\) 。

事实上，若 \(u\) 在 C 上为零，则 \(u\) 在 C - C 上为零，从而在 E 上为零，因此在 \(\mathrm{E}_{(\mathbf{C})}\) 上为零。

定理 4（克赖因–鲁特曼）. --- 设 E 是 R 上的一个完备可赋范空间。设 C 是 E 中一个闭的全突出凸锥，并设 \(u \in \mathcal{L}(\mathrm{E})\) 是满足 \(u(\mathrm{C}) \subset \mathrm{C}\) 的紧线性映射。若谱半径 \(\varrho(u)\) \textgreater0，则 \(\varrho(u)\) 是 u 的谱的孤立点，并且存在 u 的非零特征向量 \(x \in C\)，其特征值为 \(\varrho(u)\) 。

设 \(E_{(\mathbf{C})}\) 是 E 的复化，设 \(u_{(\mathbf{C})}\) 是 \(E_{(\mathbf{C})}\) 的由 u 经纯量扩张得到的自同态。\(u_{(\mathbf{C})}\) 的谱半径等于 \(\varrho(u)\)（I, p. 86）；以下简记为 \(\varrho\) 。由于 \(\varrho > 0\) ，\(u\) 的复谱不只由 0 组成。因此存在 \(\lambda_0 \in \mathrm{Sp}_\mathrm{s}(u_{(\mathbf{C})})\) 使得 \(|\lambda_0| = \varrho\)（III, p. 90 的命题 5）。设 \(e_0\) 是 \(u_{(\mathbf{C})}\) 的与 \(\lambda_0\) 相伴的谱投影。

预解式 \(\lambda \mapsto \mathrm{R}(u_{(\mathbf{C})},\lambda) = (\lambda -u_{(\mathbf{C})})^{-1}\) 在 \(u_{(\mathbf{C})}\) 的谱的余集上是全纯的（I, p. 24 的定理 1）。复数 \(\lambda_0\) 是预解式的一个极点，其留数是谱投影 \(e_0\)（I, p. 131）。

设 \(y\) 和 \(y'\) 是 E 中使得 \(y + iy' \in \mathrm{E}_{(\mathbf{C})}\) 是 \(u_{(\mathbf{C})}\) 的属于特征值 \(\lambda_0\) 的特征向量的元素。由于 \(e_0(y + iy') = y + iy' \neq 0\) ，存在元素 \(x_0 \in \mathrm{C}\) 使得 \(e_0(x_0) \neq 0\)（引理 4）。由于 C 是闭的且是尖的，其极集 \(\mathrm{C}^\circ\) 是全的（EVT, II, p. 48, 推论 1），因此存在线性形式 \(\ell_0 \in \mathrm{C}^\circ\) 使得 \(\langle e_0(x_0), \ell_0 \rangle \neq 0\) 。

考虑函数 \(f\)，它由 \(f(0)=0\) 以及

\[
f (z) = \langle \mathrm{R} (u _ {(\mathbf {C})}, z ^ {- 1}) x _ {0}, \ell_ {0} \rangle
\]

（其中 \(z \in \mathbf{C}\) 满足 \(z^{-1} \notin \mathrm{Sp}(u_{(\mathbf{C})})\)）定义。这个函数满足

\[
f (z) = \ell_ {0} \bigg (\Bigl (\sum_ {n = 0} ^ {+ \infty} z ^ {n + 1} u _ {(\mathbf {C})} ^ {n} \Bigr) x _ {0} \bigg) = \sum_ {n = 0} ^ {\infty} \langle u ^ {n} (x _ {0}), \ell_ {0} \rangle z ^ {n + 1}\tag{1}
\]

对 \(|z| < 1 / \varrho\) 成立（I, p. 24, 定理 1, d)），因而在以 0 为心、以 \(1 / \varrho\) 为半径的圆盘内是全纯的。

存在一个全纯函数 \(\widetilde{\mathrm{R}}\)，它定义在 \(\lambda_0\) 的一个开邻域 U 上，取值于 \(\mathcal{L}(\mathrm{E})\) 中，使得只要 \(z\) 属于 \(\mathrm{U} - \{\lambda_0\}\) ，就有

\[
\mathrm{R} (u _ {(\mathbf {C})}, z) = \widetilde {\mathrm{R}} (z) + \sum_ {n = 0} ^ {+ \infty} (z - \lambda_ {0}) ^ {- n - 1} (u _ {(\mathbf {C})} - \lambda_ {0}) ^ {n} e _ {0}
\]

（I, p. 83 的命题 17）。于是，对 \(z\) 满足 \(z^{-1} \in \mathrm{U}\) 且 \(z \neq 1 / \lambda_0\) 的情形，有

\[
f (z) = \langle \widetilde {\mathrm{R}} (z ^ {- 1}) x _ {0}, \ell_ {0} \rangle + \sum_ {n = 0} ^ {+ \infty} (z ^ {- 1} - \lambda_ {0}) ^ {- n - 1} \langle (u _ {(\mathbf {C})} - \lambda_ {0}) ^ {n} e _ {0} (x _ {0}), \ell_ {0} \rangle .
\]

该级数中对应于 \(n=0\) 的项是 \((z^{-1}-\lambda_{0})^{-1}\langle e_{0}(x_{0}),\ell_{0}\rangle\) 。由于 \(\langle e_{0}(x_{0}),\ell_{0}\rangle\neq0\) ，洛朗展开的唯一性（VAR, R1, p. 30, 3.3.9）表明 \(f\) 不能延拓为 \(1/\lambda_{0}\) 的一个邻域上的全纯函数。特别地，函数 \(f\) 在点 0 处的幂级数展开 (1) 的收敛半径等于 \(1/\varrho\) 。

幂级数 (1) 的系数是 \(\langle u^n(x_0),\ell_0\rangle\geqslant 0\)，因为 \(u(\mathrm{C})\subset\mathrm{C}\) 且 \(\ell_{0}\in\mathrm{C}^{\circ}\) 。由引理 3，函数 \(f\) 不能延拓为 \(1/\varrho\) 的一个邻域上的全纯函数。因此 \(u_{(\mathbf{C})}\) 的预解式不能延拓为 \(\varrho\) 的一个邻域上的全纯函数，即 \(\varrho\in\mathrm{Sp}(u)\) 。这蕴含 \(\varrho\) 是 \(u_{(\mathbf{C})}\) 的特征值（III, p. 90 的命题 5）。由于 \(\varrho\) 是实数，它也是 \(u\) 的特征值。

设 \(d \geqslant 1\) 是 \(u_{(\mathbf{C})}\) 的预解式在 \(\varrho\) 处极点的阶。设 \(e\) 是自同态

\[
e = \lim _ {z \to \varrho} (z - \varrho) ^ {d} \mathrm{R} (u _ {(\mathbf {C})}, z).
\]

它是非零的且与 \(u_{(\mathbf{C})}\) 交换，其像包含在 \(u_{(\mathbf{C})}\) 对应于 \(\varrho\) 的特征空间中（参见 I, p. 131, no. 3）。现在设 \(\ell\) 是 C° 的一个元素，\(x\) 是 C 的一个元素。由 I, p. 24, 定理 1, d)，我们有

\[
\langle e(x),\ell \rangle = \lim_{\substack{z\to \varrho \\ z > \varrho}}(z - \varrho)^{d}\sum_{n\geqslant 0}\langle u^{n}(x),\ell \rangle z^{-n - 1}\geqslant 0.
\]

双极定理（EVT, II, p. 48, 定理 1）表明 \(e(x) \in \mathrm{C}\) 对所有 \(x \in \mathrm{C}\) 成立。由于 \(\mathrm{C}\) 是全的，存在 \(x \in \mathrm{C}\) 使得 \(e(x) \neq 0\)（引理 4），于是 \(e(x)\) 属于 \(\mathrm{C}\) 且是 \(u\) 的属于特征值 \(\varrho\) 的特征向量，这正是所要证明的。

推论. --- 设 E 是 R 上的一个完备可赋范空间。设 C 是 E 中一个闭的、全的突出凸锥，并设 \(u \in \mathcal{L}(E)\) 是满足 \(u(\mathrm{C}) \subset \mathrm{C}\) 的紧线性映射。若 u 的谱半径 \(\varrho(u)\) \textgreater{} 0，则 \(\varrho(^{t}u) = \varrho(u)\) 是 \(^{t}u\) 的特征值，并且 \(^{t}u\) 有一个属于 \(\varrho(u)\) 的特征向量位于 \(C^{\circ}\) 中。

我们有 \(\varrho(^{t}u) = \varrho(u)\)（I, p. 131 的命题 3）。由 III, p. 9 的推论 1，E' 的自同态 \(^{t}u\) 是紧的。此外，我们有 \(^{t}u(C^{\circ}) \subset C^{\circ}\) ；于是结论由应用于 \(^{t}u\) 以及闭凸锥 \(C^{\circ}\) 的克赖因-鲁特曼定理得出，因为后者是突出的（由于 C 是全的）且是全的（由于 C 是突出的）。

注. --- 在克赖因–鲁特曼定理中，假设 \(\varrho(u) > 0\) 一般不能省略。例如，设 V 是巴拿赫空间 \(\mathcal{C}([0,1])\) 上由

\[
\mathrm{V} (f) (x) = \int_ {0} ^ {x} f (y) d y
\]

（其中 \(f \in \mathcal{C}([0,1])\)）所定义的自同态。映射 V 是紧的，且其谱半径为零（I, p. 187 的习题 1）。它保持 \(\mathcal{C}([0,1])\) 中由正函数组成的闭的、全的、突出的凸锥，并且没有特征值（出处同上）。

下面的命题刻画了保持一个锥的自同态具有严格正谱半径的一个充分条件，并进而细化了克赖因-鲁特曼定理。

命题 8. --- 设 E 是 R 上的一个非零完备可赋范空间。设 C 是 E 中内部 \(\mathring{\mathrm{C}}\) 非空的闭突出凸锥，并设 \(u \in \mathcal{L}(\mathrm{E})\) 是满足 \(u(\mathrm{C} - \{0\}) \subset \mathring{\mathrm{C}}\) 的紧线性映射。

\begin{enumerate}
\def\labelenumi{\alph{enumi})}
\item
  有 \(\varrho(u) > 0\)，并且存在一个向量 \(x_0\)，它是 \(u\) 的位于 \(\mathring{C}\) 中、属于特征值 \(\varrho(u)\) 的特征向量；
\item
  u 的对应于 \(\varrho(u)\) 的谱投影的秩为 1；
\item
  设 \(\lambda \neq \varrho(u)\) 是 \(u_{(\mathbf{C})}\) 的任意一个特征值，则 \(|\lambda| < \varrho(u)\)；
\item
  设 F 是 E 的在 u 下稳定且使 \((\mathrm{C}-\{0\})\cap\mathrm{F}\) 非空的子空间。则 \(x_{0}\in F\) 。特别地，u 在 C 中的特征向量只有 \(x_{0}\) 的倍向量。
\end{enumerate}

我们先证明两个预备引理。

引理 5. --- 设 E 是实巴拿赫空间，C 是 E 中的一个凸锥，\(u \in \mathcal{L}(\mathrm{E})\) 满足 \(u(\mathrm{C} - \{0\}) \subset \mathring{\mathrm{C}}\) 。设 \(\ell \in \mathrm{C}^{\circ}\) 是 \(^{t}u\) 的一个特征向量。则 \(\ell\) 的核在 \(u\) 下稳定且与 \(\mathrm{C} - \{0\}\) 不相交。

设 \(\lambda \in \mathbf{R}\) 满足 \(^t u(\ell) = \lambda \ell\) 。对每个 \(x \in \mathrm{Ker}(\ell)\)，我们有

\[
\langle u (x), \ell \rangle = \langle x, ^ {t} u (\ell) \rangle = \lambda \langle x, \ell \rangle = 0,
\]

因而 \(\operatorname{Ker}(\ell)\) 在 \(u\) 下稳定。

设 \(x\) 是 \(\operatorname{Ker}(\ell)\) 的一个非零元素。对 E 的每个元素 \(y\)，只要 \(\langle y, \ell \rangle < 0\)，就有 \(\langle u(x) + y, \ell \rangle < 0\)，由此推出 \(u(x) + y \notin \mathrm{C}\)，因为 \(\ell \in \mathrm{C}^\circ\) 。我们推得 \(u(x)\) 不属于 \(\mathring{\mathrm{C}}\) 。由于 \(u(\mathrm{C} - \{0\}) \subset \mathring{\mathrm{C}}\) ，这蕴含 \(x \notin \mathrm{C}\) 。

引理 6. --- 设 E 是有限维实向量空间，B 是 E 中 0 的一个紧凸邻域。设 u 是 E 的满足 \(u(B) \subset \mathring{B}\) 的自同态。则 u 的复谱包含于 C 的单位圆盘，特别地，它与 C 中以 0 为心、以 1 为半径的圆不相交。

将 B 换成 \(B \cap (-B)\) ，可化归为 B 是均衡集的情形。此时 B 的度规 p（EVT, II, p. 28, 习题 3）是 E 上的一个范数，并且它给出 E 的拓扑（EVT, I, p. 14, 定理 2）。假设蕴含 \(p(u(x)) < p(x)\) 对 E - \{0\} 的每个元素 x 成立，从而 u 的谱半径 \textless{} 1；由于它也是 \(u_{(\mathbf{C})}\) 的谱半径（参见 I, p. 86），结论得证。

现在来证明该命题。

用 \(x \preccurlyeq y\) 表示 E 上与凸锥 C 相伴的序关系（EVT, II, p. 13, n° 5），即 \(x \preccurlyeq y\) 当且仅当 \(y - x \in \mathrm{C}\) 。我们有 \(u(x) \preccurlyeq u(y)\)，只要 \(x \preccurlyeq y\)。由于 \(\mathring{\mathrm{C}}\) 非空，由 C 生成的向量空间 C -- C（EVT, II, p. 12, 推论 1）包含 0 的一个邻域，因此锥 C 是全的。

我们来证明断言 \(a\) )。设 \(\varrho = \varrho(u)\) 。设 \(y_0\) 是 \(\mathring{\mathrm{C}}\) 的一个元素。我们有 \(y_0 \neq 0\) 。取 \(r > 0\) 使得以 \(y_0\) 为心、以 \(r\) 为半径的闭球包含于 C。对每个 \(y \in \mathrm{E} - \{0\}\) ，有 \(y_0 - r \| y\|^{-1}y \in \mathrm{C}\) ，因而 \(y \preccurlyeq r^{-1}\| y\| y_0\) 对所有 \(y \in \mathrm{E}\) 成立。

由于 \(y_{0} \neq 0\)，假设蕴含 \(u(y_{0}) \in \mathring{\mathrm{C}}\) 。因此存在 t \textgreater{} 0 使得 \(tu(y_{0}) - y_{0} \in \mathrm{C}\) 。令 v = tu。这是 E 的一个紧自同态，满足 \(v(\mathrm{C}) \subset \mathrm{C}\) 且 \(v(y_{0}) \succcurlyeq y_{0}\) 。对每个整数 \(n \geqslant 1\)，我们有

\[
y _ {0} \preccurlyeq v ^ {n} (y _ {0}) \preccurlyeq r ^ {- 1} \| v ^ {n} (y _ {0}) \| y _ {0} \preccurlyeq r ^ {- 1} \| v ^ {n} \| \| y _ {0} \| y _ {0},
\]

因而 \((r^{-1}\|v^{n}\|\|y_{0}\|-1)y_{0}\in\mathrm{C}\) 。由于 C 是突出的，这蕴含 \(t^{n}\|u^{n}\|=\|v^{n}\|\geqslant r/\|y_{0}\|\) ，从而 \(\varrho\geqslant t^{-1}>0\)（I, p. 20, 命题 1）。

由克赖因–鲁特曼定理（定理 4），实数 \(\varrho\) 是 \(u\) 的特征值，且存在一个向量 \(x_0\)，它是 \(u\) 的位于 C 中、属于特征值 \(\varrho\) 的特征向量。由假设有 \(\varrho x_0 = u(x_0) \in \mathring{\mathrm{C}}\)，因而 \(x_0 \in \mathring{\mathrm{C}}\) 。这就证明了断言 \(a\) )。

设 K 是 \(u_{(\mathbf{C})}\) 的谱与 C 中以 0 为心、以 \(\varrho\) 为半径的圆的交。由于 u 是紧的且 \(\varrho > 0\)，集合 K 是有限的，并且谱投影 \(e_{K}\) 的像是 \(\mathrm{E}_{(\mathbf{C})}\) 的一个有限维子空间（III, p. 83 的命题 2 与 III, p. 90 的命题 5）。由于 K 在复共轭下稳定，\(e_{K}\) 的像是 \(\mathrm{E}_{(\mathbf{C})}\) 中在 R 上有理的有限维子空间（A, V, p.60, 命题 6）；用 F 表示 E 中使 \(\mathrm{F}_{(\mathbf{C})}\) 等于 \(e_{K}\) 的像的子空间。空间 F 在 u 下稳定且包含 u 对应于 \(\varrho\) 的特征空间，因此 F 非零。

为证明断言 \(b)\) 与 \(c),\)，只需证明 \(F\) 是一维的。

用 \(v\) 表示由 \(u\) 通过过渡到子空间得到的 F 的自同态。设 \(C_F = C \cap F\) 。这是 F 中内部非空的闭突出凸锥（因为 \(x_0 \in \mathring{C} \cap F\)），因而是全的；它满足 \(v(C_F - \{0\}) \subset \mathring{C}_F\)，特别地在 \(v\) 下稳定。由于 \(\varrho(v) \leqslant \varrho\) 且 \(x_0 \in F\)，我们有 \(\varrho(v) = \varrho\)。由定理 4 的推论（应用于 \(C_F\) 和 \(v\)），存在一个向量 \(\ell\)，它是 \(^t v\) 的位于 \(C_F^\circ\) 中、属于特征值 \(\varrho(v) = \varrho > 0\) 的特征向量。

F 的子空间 \(\mathrm{Ker}(\ell)\) 在 v 下稳定。设 w 是 \(\mathrm{Ker}(\ell)\) 的、由 \(\varrho^{-1}v\) 通过过渡到商子空间而得到的自同态；它的谱包含于单位圆。集合 B 由这样的 \(y \in \mathrm{Ker}(\ell)\) 组成：\(x_{0} + y \in C_{F}\)；它是 \(\mathrm{Ker}(\ell)\) 中 0 的一个闭凸邻域。对每个 \(y \in B\) 和每个 \(z \in \mathrm{Ker}(\ell)\)，我们有

\[
x _ {0} + (w (y) + z) = \varrho^ {- 1} (v (x _ {0} + y) + \varrho z)
\]

当 z 的范数充分小时它属于 \(C_{F}\)，因而 \(w(y) \in \mathring{\mathrm{B}}\) 。集合 B 是有界的：事实上，若存在一个序列 \((y_{n})_{n \in \mathbf{N}}\)，它含于 \(\operatorname{Ker}(\ell)\)，使得 \(\|y_{n}\| \to +\infty\) 且 \(x_{0} + y_{n} \in C_{F}\)，那么在过渡到子列后，将有 \(y_{n}/\|y_{n}\| \to y\)，其中 \(y \in \operatorname{Ker}(\ell)\) 非零，而 \(y = \lim \|y_{n}\|^{-1}(x_{0} + y_{n})\) 将属于 \(C_{F}\)，这与引理 5 矛盾。因此集合 B 是紧的。

于是，由应用于 \(\mathrm{Ker}(\ell)\)、B 和 \(w\) 的引理 6 推出，\(w\) 的复谱与以 0 为心、以 1 为半径的圆不相交；这蕴含 \(w\) 的谱是空的，即 \(\mathrm{Ker}(\ell)\) 化为 \(\{0\}\)，从而 F 的维数为 1。因此断言 \(b\) 和 \(c\) 得证。

线性映射 \({}^t u\) 是紧的（III, p. 9 的推论 1），且 \({}^t u(\mathrm{C}^\circ) \subset \mathrm{C}^\circ\)；此外 \(\varrho(^t u) = \varrho > 0\)（I, p. 131 的命题 3）。由定理 4 的推论，在 \(\mathrm{C}^\circ\) 中存在一个向量 \(\ell \neq 0\)，它是 \({}^t u\) 的属于特征值 \(\varrho\) 的特征向量。\(\ell\) 的核在 \(u\) 下稳定且与 \(\mathrm{C} - \{0\}\) 不相交（引理 5）。特别地，有 \(x_0 \notin \operatorname{Ker}(\ell)\) 。

设 F 是 E 的一个在 u 下稳定的子空间。我们有分解 \(\mathrm{F} = (\mathrm{F} \cap \mathbf{R}x_{0}) \oplus (\mathrm{F} \cap \operatorname{Ker}(\ell))\) 。若 F 不包含 \(x_{0}\)，则 \(\mathrm{F} \subset \operatorname{Ker}(\ell)\)，从而 F 与 \(C - \{0\}\) 不相交。这就证明了 d)，命题证毕。

推论. --- 设 E 是 R 上的一个非零完备可赋范空间。设 C 是 E 中内部 \(\dot{C}\) 非空的闭凸突出锥，并设 u 是 E 的满足 \(u(\mathrm{C}) \subset \mathrm{C}\) 的紧自同态。假设存在整数 \(k \geqslant 1\) 使得 \(u^{k}(\mathrm{C} - \{0\}) \subset \mathring{\mathrm{C}}\) 。

\begin{enumerate}
\def\labelenumi{\alph{enumi})}
\item
  有 \(\varrho(u) > 0\)，并且存在一个向量 \(x_0\)，它是 \(u\) 的位于 \(\dot{C}\) 中、属于特征值 \(\varrho(u)\) 的特征向量；
\item
  u 的对应于 \(\varrho(u)\) 的谱投影的秩为 1；
\item
  设 \(\lambda \neq \varrho(u)\) 是 \(u_{(\mathbf{C})}\) 的任意一个特征值，则 \(|\lambda| < \varrho(u)\)；
\item
  u 在 C 中的特征向量只有 \(x_0\) 的倍向量。
\end{enumerate}

设 \(\varrho = \varrho(u)\) 。可以把定理 4 应用于 \(u\)，把命题 8 应用于 \(u^k\) 。因此存在一个向量 \(x_0\)，它是 \(u\) 的位于 C 中、属于特征值 \(\varrho\) 的特征向量。由于 \(0 < \varrho(u^k) = \varrho^k\)（I, p. 21 的公式 (4)），特别地有 \(\varrho > 0\)；又由于 \(x_0\) 是 \(u^k\) 的属于特征值 \(\varrho(u^k)\) 的特征向量，我们有 \(x_0 \in \mathring{\mathrm{C}}\) 。

我们有 \(\mathrm{Sp}(u_{(\mathbf{C})}^{k})=\mathrm{Sp}(u_{(\mathbf{C})})^{k}\)（I, p. 2 的注 4），因此每个特征值 \(\lambda\in C\)（属于 \(u_{(\mathbf{C})}\)），只要 \(|\lambda|=\varrho\)，就满足 \(\lambda^{k}=\varrho^{k}\)；由于对应于 \(\lambda\) 的每个特征向量都是 \(u^{k}\) 的属于 \(\varrho^{k}\) 的特征向量，从而与 \(x_{0}\) 成比例，故有 \(\lambda=\varrho\) 。

若 u 的对应于特征值 \(\varrho\) 的谱投影的秩至少为 2，则 \(u^{k}\) 的对应于特征值 \(\varrho^{k}\) 的谱投影的秩也至少为 2（参见 I, p. 129, no. 2）。

最后，若 \(x \in C\) 是 u 的一个特征向量，则它也是 \(u^{k}\) 的特征向量，从而与 \(x_{0}\) 成比例。

设 n 是一个 \(\geqslant 1\) 的整数。集合 \(R_{+}^{n}\) 是 \(R^{n}\) 中的闭凸尖突出锥，其非空内部等于 \((\mathbf{R}_{+}^{*})^{n}\) 。

设 \(A = (a_{i,j})\) 是一个 \((n,n)\) 型实矩阵，\(u_A\) 是由 \(x \mapsto Ax\) 给出的 \(\mathbf{R}^n\) 的自同态。设 \(\varrho = \varrho(u_A)\) 是它的谱半径。

引理 7. --- 有 \(u_A(\mathbf{R}_+^n) \subset \mathbf{R}_+^n\)，当且仅当对所有 i 和 j 有 \(a_{i,j} \geqslant 0\)；有 \(u_A(\mathbf{R}_+^n - \{0\}) \subset (\mathbf{R}_+^*)^n\)，当且仅当对所有 i 和 j 有 \(a_{i,j} > 0\)。

设 \((e_{1},\ldots,e_{n})\) 是 \(R^{n}\) 的典范基。向量 \(e_{i}\) 属于 \(R_{+}^{n}\)，且对所有 i 有 \(u_{A}(e_{i})\in\mathbf{R}_{+}^{n}\)（相应地，对所有 i 有 \(u_{A}(e_{i})\in(\mathbf{R}_{+}^{*})^{n}\)），当且仅当对所有 i 和 j 有 \(a_{i,j}\geqslant0\)（相应地，对所有 i 和 j 有 \(a_{i,j}>0\)）。由线性性即得结论。

定理 5（佩龙–弗罗贝尼乌斯）. --- a) 若 \(a_{i,j} \geqslant 0\)，其中 \(i\) 和 \(j\) 取遍 \(\{1,\ldots,n\}\)，则实数 \(\varrho\) 是 \(A\) 的特征值；b) 若 \(a_{i,j} > 0\)，其中 \(i\) 和 \(j\) 取遍 \(\{1,\ldots,n\}\)，则 \(\varrho > 0\)，且 \(u_A\) 对应于 \(\varrho\) 的准素空间，即 \(A - \varrho 1_{\mathbf{R}^n}\) 的幂零空间，维数为 1，并由一个向量 \(x_0 \in (\mathbf{R}_+^*)^n\) 生成。\(A\) 没有其他在 \(\mathbf{R}_+^n\) 中有特征向量的特征值，并且所有复特征值 \(\lambda \neq \varrho\)（\(A\) 的）都满足 \(|\lambda| < \varrho\) 。

若 \(\varrho = 0\)，断言 \(a\) ) 是初等的，因为此时 0 是 \(A\) 的特征值。若 \(\varrho > 0\)，则由引理 7，断言 \(a\) ) 可由定理 4 得出。出于同样的理由，断言 \(b\) ) 是命题 8 的推论。

注. --- 假设存在整数 \(k \geqslant 1\) 使得 \(A^{k}\) 的所有系数都 \textgreater0（这样的矩阵有时称为本原矩阵）。那么，由命题 8 的推论，谱半径 \(\varrho\) 是 A 的特征值，\(u_{A}\) 对应于 \(\varrho\) 的准素空间维数为 1，且由一个向量 \(x_{0} \in (\mathbf{R}_{+}^{*})^{n}\) 生成。此外，A 没有其他在 \(R_{+}^{n}\) 中有特征向量的复特征值，并且 A 的所有复特征值 \(\lambda \neq \varrho\) 满足 \(|\lambda| < \varrho\) 。

\specialchapter{习题}

\exercisesection{}

仅在习题 1 至 5 中，K 表示一个完备的非阿基米德、非离散的赋值域，其赋值记为 \(x \mapsto |x|\) 。我们还用 \(\mathrm{G} \subset \mathbf{R}_{+}^{*}\) 表示 \(\mathrm{K}^{*}\) 在映射 \(x \mapsto |x|\) 下的像。所考虑的巴拿赫空间都是 \(\mathrm{K}\) 上的巴拿赫空间。

\begin{enumerate}
\def\labelenumi{\arabic{enumi})}
\tightlist
\item
  \begin{enumerate}
  \def\labelenumii{\alph{enumii})}
  \tightlist
  \item
    设 E 是 K 上的一个巴拿赫空间。E 上存在一个范数 \(x \mapsto \|x\|\)，它给出 E 的拓扑，且使得 \(\|x\| \in \overline{G}\) 对所有 \(x \in E\) 成立。这样的范数称为圆范数。若 K 的赋值不是离散的，则每个给出 E 的拓扑的范数都是圆范数。
  \end{enumerate}
\end{enumerate}

\begin{enumerate}
\def\labelenumi{\alph{enumi})}
\setcounter{enumi}{1}
\item
  设 E 和 F 是 K 上的巴拿赫空间。我们在空间 \(\mathcal{L}(\mathrm{E};\mathrm{F})\) 上赋予范数 \(\| u\| = \sup_{x\in \mathrm{E}- \{0\}}\| u(x)\| /\| x\|\) 。空间 \(\mathcal{L}(\mathrm{E};\mathrm{F})\) 是 K 上的巴拿赫空间。若 E 上的范数是圆范数，则 \(\| u\| = \sup_{\| x\| \leqslant 1}\| u(x)\|\) 。若 F 上的范数也是圆范数，则 \(\mathcal{L}(\mathrm{E};\mathrm{F})\) 上的范数是圆范数。我们记 \(\mathrm{E}' = \mathcal{L}(\mathrm{E};\mathrm{K})\)，并称 \(\mathrm{E}'\) 为 E 的对偶空间。
\item
  设 I 是一个集合，F 是 K 上的一个巴拿赫空间，\(\mathcal{C}_{0}(\mathrm{I},\mathrm{F})\) 是由函数 \(f\colon I\to F\) 组成的 K-向量空间，其中 \(\|f(i)\|\) 沿 I 的有限子集的余集滤子趋于 0。赋以范数 \(\|f\|=\sup_{i\in I}\|f(i)\|\) 后，空间 \(\mathcal{C}_{0}(\mathrm{I},\mathrm{F})\) 是 K 上的巴拿赫空间，且若 F 的范数是圆范数，则它的范数也是圆范数。我们记 \(\mathcal{C}_{0}(\mathrm{I})=\mathcal{C}_{0}(\mathrm{I},\mathrm{K})\) 。
\end{enumerate}

\(d\) ) 设 I 是一个集合。对每个 \(i\in \mathrm{I}\) ，用 \(\varphi_{i}\in \mathcal{C}_{0}(\mathrm{I})\) 表示 \(\{i\}\) 的特征函数。设 F 是 K 上的一个巴拿赫空间。映射 \(u \mapsto (u(\varphi_i))_{i \in \mathrm{I}}\) 是从巴拿赫空间 \(\mathcal{L}(\mathcal{C}_0(\mathrm{I});\mathrm{F})\) 到巴拿赫空间 \(\mathcal{B}(\mathrm{I};\mathrm{K})\) 的同构（EVT, I, p. 4, 例；参见 EVT, I, p. 23, 习题 7）。

\begin{enumerate}
\def\labelenumi{\alph{enumi})}
\setcounter{enumi}{4}
\item
  假设 K 的赋值是离散的。设 E 是 K 上的范数为圆范数的巴拿赫空间。存在一个集合 I，使得 E 等距同构于巴拿赫空间 \(\mathcal{C}_{0}(\mathrm{I})\)（参见 EVT, I, p. 23, 习题 7）。
\item
  假设 K 的赋值是离散的。设 E 是赋有圆范数的 K 上的巴拿赫空间，F 是 E 的一个闭向量子空间。存在 E 的一个连续投影，其像为 F 且范数 \(\leqslant 1\) 。
\end{enumerate}

\begin{enumerate}
\def\labelenumi{\arabic{enumi})}
\setcounter{enumi}{1}
\tightlist
\item
  设 E 和 F 是 K 上的巴拿赫空间。我们称 \(u \in \mathcal{L}(\mathrm{E};\mathrm{F})\) 是完全连续的，如果它属于 \(\mathcal{L}(\mathrm{E};\mathrm{F})\) 中由有限秩线性映射组成的空间 \(\mathcal{L}^{\mathrm{f}}(\mathrm{E};\mathrm{F})\) 的闭包。我们用 \(\mathcal{L}^{\mathrm{c}}(\mathrm{E};\mathrm{F})\) 表示从 E 到 F 的完全连续线性映射组成的空间，并令 \(\mathcal{L}^{\mathrm{c}}(\mathrm{E}) = \mathcal{L}^{\mathrm{c}}(\mathrm{E};\mathrm{E})\) 。 a) 空间 \(\mathcal{L}^{\mathrm{c}}(\mathrm{E})\) 是 \(\mathcal{L}(\mathrm{E})\) 的一个闭双边理想。设 \(u \in \mathcal{L}(\mathrm{E};\mathrm{F})\) 。映射 \(\ell \mapsto \ell \circ u\) 是从 \(F'\) 到 \(E'\) 的连续线性映射，记为 \({}^{t}u\) 。若 \(u \in \mathcal{L}^{\mathrm{c}}(\mathrm{E};\mathrm{F})\) ，则 \({}^{t}u \in \mathcal{L}^{\mathrm{c}}(\mathrm{F}'; \mathrm{E}')\) 。
\end{enumerate}

\begin{enumerate}
\def\labelenumi{\alph{enumi})}
\setcounter{enumi}{1}
\item
  设 J 是一个集合，\(\mathrm{F} = \mathcal{C}_{0}(\mathrm{J})\) 。设 \(u \in \mathcal{L}(\mathrm{E};\mathrm{F})\) 。对每个 \(j \in J\)，映射 \(\ell_{j} \colon x \mapsto u(x)(j)\) 是 \(E'\) 的一个元素。我们有 \(\|u\| = \sup_{j \in J} \| \ell_{j} \|\)，并且映射 \(u \mapsto (\ell_{j})_{j \in J}\) 通过过渡到子空间诱导一个从 \(\mathcal{L}^{\mathrm{c}}(\mathrm{E};\mathrm{F})\) 到巴拿赫空间 \(\mathcal{C}_{0}(\mathrm{J};\mathrm{E}')\) 的同构。
\item
  若 K 是局部紧的，则 \(u \in \mathcal{L}(\mathrm{E}; \mathrm{F})\) 是完全连续的，当且仅当 E 的每个有界子集在 u 下的像在 F 中是相对紧的。（可将此结果与 III, p. 16 的命题 11 相比较。）
\end{enumerate}

\begin{enumerate}
\def\labelenumi{\arabic{enumi})}
\setcounter{enumi}{2}
\tightlist
\item
  我们用 A 表示 K 的由元素 \(x \in \mathrm{K}\)（满足 \(|x| \leqslant 1\)）组成的子环；它是一个局部环，其极大理想由 \(x \in \mathrm{A}\)（满足 \(|x| < 1\)）组成。
\end{enumerate}

设 I 是一个集合，E 是巴拿赫空间 \(\mathcal{C}_{0}(\mathrm{I})\)（参见习题 1 的问题 c)）。对 \(i \in I\)，用 \(\varphi_{i} \in E\) 表示 \(\{i\}\) 的特征函数。设 \(u \in \mathcal{L}^{\mathrm{c}}(\mathrm{E})\) 满足 \(\|u\| \leqslant 1\) 。

\begin{enumerate}
\def\labelenumi{\alph{enumi})}
\item
  子空间 \(E_{0}\) 由这样的元素 \(x \in E\) 组成：\(\|x\| \leqslant 1\)；它在 u 下稳定。
\item
  对 A 的每个非零理想 a，线性映射 u 通过过渡到商定义一个自同态 \(u_{a}\)，它作用于 A/a-模 \(E_{a} = E_{0}/aE_{0}\) 。这个 A/a-模是自由的，且 \(u_{a}\) 的像包含于一个有限生成的子模中。
\item
  存在唯一的形式级数 \(f_{u} \in A[[t]]\)，使得对 A 的每个非零理想 a，\(f_{u}\) 在 \((A/\mathfrak{a})[[t]]\) 中的像等于 A, VIII, p. 455 中定义的多项式 \(\det(1 - tu_{\mathfrak{a}})\)。
\end{enumerate}

\(d)\) 设 \(v \in \mathcal{L}^{\mathrm{c}}(\mathrm{E})\)，并设 \(c \in \mathrm{K}^{*}\) 满足 \(\|cv\| \leqslant 1\) 。形式级数 \(f_{cv}\) 与满足此性质的 \(c\) 的选取无关。

我们将用 \(\det(1-tv)\) 表示形式级数 \(f_v\)，这里 \(c \in \mathbf{K}^*\) 为任意满足 \(\|cv\| \leqslant 1\) 的元素。

¶ 4) 我们沿用上一习题的记号。

\begin{enumerate}
\def\labelenumi{\alph{enumi})}
\item
  设 \(v \in \mathcal{L}^c(\mathrm{E})\) 。形式级数 \(\det(1 - tv)\) 的收敛半径是无穷大。
\item
  设 \((v_{n})_{n\in\mathbf{N}}\) 是 \(\mathcal{L}^{c}(\mathrm{E})\) 中收敛于 \(v\in\mathcal{L}^{c}(\mathrm{E})\) 的一个序列。形式级数序列 \(\det(1-tv_{n})\) 对于形式级数系数的逐点收敛拓扑收敛于 \(\det(1-tv)\)。
\item
  设 \(v \in \mathcal{L}^{c}(E)\) 。对 K 的每个扩张 L 和每个 \(x \in L\)，我们用 \(\det(1 - xv)\) 表示形式幂级数 \(\det(1 - tv)\) 在 x 处的值，它由 a) 良定义。我们有 \(\det(1 + u)\det(1 + v) = \det(1 + u + v + u \circ v)\) 。
\item
  设 J 是一个集合，\(\mathrm{F} = \mathcal{C}_{0}(\mathrm{J})\) 。对所有 \(u \in \mathcal{L}^{\mathrm{c}}(\mathrm{E};\mathrm{F})\) 和 \(v \in \mathcal{L}^{\mathrm{c}}(\mathrm{F};\mathrm{E})\)，作为形式幂级数有 \(\det(1 - tu \circ v) = \det(1 - tv \circ u)\) 。
\item
  设 \(v \in \mathcal{L}^{\mathrm{c}}(\mathrm{E})\) 。我们用 \(\operatorname{Tr}(v)\) 表示 \(\det(1 - tu)\) 中 t 的系数的相反数。若 K 的特征为零，则
\end{enumerate}

\[
\det (1 - t u) = \exp \left(- \sum_ {m \geqslant 1} \frac {\operatorname{Tr} \left(u ^ {m}\right)}{m} t ^ {m}\right)
\]

作为形式幂级数成立。

¶ 5) 我们保留上一习题的记号。设 \(v \in \mathcal{L}^{\mathrm{c}}(\mathrm{E})\) 。

\begin{enumerate}
\def\labelenumi{\alph{enumi})}
\item
  形式幂级数 \(\det(1 - tv)\sum_{k\geqslant 0}t^{k}v^{k}\)（\(\mathcal{L}(\mathrm{E})[[t]]\) 中的）的收敛半径为无穷大。
\item
  设 \(\lambda \in \mathrm{K}\) 。E 的自同态 \(1_{\mathrm{E}} - \lambda v\) 可逆，当且仅当 \(\det(1 - \lambda v) \neq 0\) 。
\item
  设 \(\lambda \in K\) 满足 \(\det(1 - \lambda v) = 0\) 。存在 E 的唯一的闭子空间 \(I_{\lambda}\) 和 \(N_{\lambda}\)，它们在 \(v\) 下稳定，使得
\end{enumerate}

\begin{enumerate}
\def\labelenumi{(\roman{enumi})}
\item
  空间 E 是 \(I_{\lambda}\) 与 \(N_{\lambda}\) 的拓扑直和；
\item
  \(I_{\lambda}\) 的由 \(1_{E} - \lambda v\) 通过过渡到子空间得到的自同态是可逆的；
\item
  \(N_{\lambda}\) 的由 \(1_{E} - \lambda v\) 通过过渡到子空间得到的自同态是幂零的。
\end{enumerate}

\(d)\) 设 \(h \geqslant 1\) 是 det(1 - tv) 的零点 \(\lambda\) 的阶。我们有 \(\dim(N_{\lambda}) = h\) 。

应将本习题的结果与 III, p. 77 的 no. 5 中关于 R 或 C 上巴拿赫空间的结果相比较。

在下文中，K 表示域 R 或 C，且所有拓扑向量空间都是 K 上的。

\begin{enumerate}
\def\labelenumi{\arabic{enumi})}
\setcounter{enumi}{5}
\tightlist
\item
  设 E 是空间 \(C^{N}\)，赋以与半范数序列 \((p_{n})_{n\in\mathbf{N}}\) 相伴的拓扑，其中 \(p_{n}(x)=|x_{n}|\) 对 E 的每个元素 \(x=(x_{n})_{n\in\mathbf{N}}\) 成立。
\end{enumerate}

\begin{enumerate}
\def\labelenumi{\alph{enumi})}
\item
  空间 E 是一个弗雷歇空间。
\item
  E 的紧自同态组成的集合 \(\mathcal{L}^{\mathrm{c}}(\mathrm{E})\) 在 \(\mathcal{L}(\mathrm{E})\) 中不是闭的。
\end{enumerate}

\begin{enumerate}
\def\labelenumi{\arabic{enumi})}
\setcounter{enumi}{6}
\item
  设 E 和 F 是局部凸拓扑向量空间。若 F 是可度量化的，则每个紧线性映射 \(E \rightarrow F\) 的像是可数型的。
\item
  构造希尔伯特空间 E 的一个非紧自同态 \(u\) 的例子，使得 \(u^2\) 是紧的。
\item
  对每一对 \((x,y)\in]0,1]^{2}\) ，存在一个整数 \(n\geqslant1\) 和一个整数 k 满足 \(0\leqslant k\leqslant2^{n-1}-1\)，使得
\end{enumerate}

\[
2 ^ {- n} <   x \leqslant 2 ^ {- n + 1}, \quad \frac {2 k}{2 ^ {n}} <   y \leqslant \frac {2 k + 2}{2 ^ {n}}.
\]

然后我们如下定义 N: ]0,1]×]0,1] → C：

\[
\mathrm{N} (x, y) = \left\{ \begin{array}{l l} 2 & \text { if } \frac {2 k}{2 ^ {n}} <   y \leqslant \frac {2 k + 1}{2 ^ {n}} \\ 0 & \text { if } \frac {2 k + 1}{2 ^ {n}} <   y \leqslant \frac {2 k + 2}{2 ^ {n}}. \end{array} \right.
\]

\begin{enumerate}
\def\labelenumi{\alph{enumi})}
\tightlist
\item
  证明公式
\end{enumerate}

\[
u (f) (x) = \int_ {[ 0, 1 ]} \mathrm{N} (x, y) f (y) d y
\]

在 \(L^{1}([0,1])\) 上定义了一个关于勒贝格测度的连续线性映射 u。

\begin{enumerate}
\def\labelenumi{\alph{enumi})}
\setcounter{enumi}{1}
\item
  证明自同态 \(u\) 不是紧的。
\item
  自同态 \(u^2\) 是紧的。
\end{enumerate}

\(d)\) 设 \((x_{n})_{n\in \mathbf{N}}\) 是 \(\mathrm{L}^1 ([0,1])\) 中的一个有界序列。序列 \((u(x_n))_{n\in \mathbf{N}}\) 有一个弱收敛的子列。

\begin{enumerate}
\def\labelenumi{\alph{enumi})}
\setcounter{enumi}{4}
\item
  设 \((x_{n})_{n\in \mathbf{N}}\) 是 \(\mathrm{L}^1 ([0,1])\) 中的一个弱收敛序列。序列 \((u(x_{n}))_{n\in \mathbf{N}}\) 在 \(\mathrm{L}^1 ([0,1])\) 中收敛。
\item
  设 E 是这样的拓扑向量空间：其元素是有界可测的函数 \(f\colon[0,1]\to\mathbf{C}\)，赋以强于 \(\mathrm{L}^{2}([0,1])\) 的拓扑以及相对于 \(\mathrm{L}^{1}([0,1])\) 的弱拓扑的最细拓扑。则由
\end{enumerate}

\[
v (f) (x) = \int_ {[ 0, 1 ]} \overline {{\mathrm{N} (y , x)}} f (y) d y
\]

定义的线性映射 v 是 E 的一个紧自同态。它的转置 \(^{t}v\) 对于 E' 上的强拓扑不是紧的。

\begin{enumerate}
\def\labelenumi{\arabic{enumi})}
\setcounter{enumi}{9}
\tightlist
\item
  设 p 和 r 是满足 \(1 \leqslant p < r\) 的实数。设 u 是从 \(\ell_{\mathrm{K}}^{r}(\mathbf{N})\) 到 \(\ell_{\mathrm{K}}^{p}(\mathbf{N})\) 的一个连续线性映射。
\end{enumerate}

我们用 \(r_{n}\)（相应地，\(p_{n}\)）表示投影

\[
(x _ {i}) _ {i \in \mathbf {N}} \mapsto (x _ {0}, \ldots , x _ {n - 1}, 0, 0, \ldots)
\]

到 \(\ell_{\mathrm{K}}^{r}(\mathbf{N})\) 上（相应地，到 \(\ell_{\mathrm{K}}^{p}(\mathbf{N})\) 上）。我们用 \(\| \cdot \| _p\)（相应地，\(\| \cdot \| _r\)）表示 \(\ell_{\mathrm{K}}^{p}(\mathbf{N})\)（相应地，\(\ell_{\mathrm{K}}^{r}(\mathbf{N})\)）上的范数。

\begin{enumerate}
\def\labelenumi{\alph{enumi})}
\item
  证明对每个 \(n \in \mathbf{N}\) ，有 \(\|p_n u(1 - r_m)\| \to 0\)（当 \(m \to +\infty\)），其中 1 表示 \(\ell_{\mathrm{K}}^{r}(\mathbf{N})\) 的恒同态射。
\item
  对 \(n\) 、 \(m \in \mathbf{N}\) ，令 \(u_{n,m} = (1 - p_n)u(1 - r_m)\) 。设 \(\alpha = (\alpha_n)_{n \in \mathbf{N}}\) 是一个严格正实数的序列，满足 \(\|\alpha\|_r = 1\) 且 \(\alpha\) 不属于 \(\ell_{\mathrm{K}}^p(\mathbf{N})\)，并设 \((\varepsilon_n)\) 是一个严格正实数的序列，使得级数 \(\sum \varepsilon_n\alpha_n\) 收敛。
\end{enumerate}

若 \(u\) 不是紧的，证明存在 \(\delta > 0\) 、整数的严格递增序列 \((n_i)_{i \in \mathbf{N}}\) 与 \((m_i)_{i \in \mathbf{N}}\)，以及两个序列：\((x_i)_{i \in \mathbf{N}}\)（由 \(\ell_{\mathrm{K}}^r(\mathbf{N})\) 的元素组成）与 \((y_i)_{i \in \mathbf{N}}\)（由 \(\ell_{\mathrm{K}}^p(\mathbf{N})\) 的元素组成），具有下列性质：

\begin{enumerate}
\def\labelenumi{(\roman{enumi})}
\item
  有 \(\|x_{i}\|_{r}=1\) 对所有 \(i\in N\) 成立；
\item
  集合 \(S_{i} = \{n \in N \mid x_{i,n} \neq 0\}\) 两两不相交，且集合 \(T_{i} = \{n \in N \mid y_{i,n} \neq 0\}\) 两两不相交；
\item
  有 \(\| y_i\| _p\geqslant \delta\)
\item
  有 \(\|u(x_{i})\|_{q} > \delta - 2\varepsilon_{i}\) 。
\end{enumerate}

（用归纳法进行证明，并利用公式

\[
u = u _ {n, m} + p _ {n} u (1 - r _ {m}) + (1 - p _ {n}) u r _ {m} + p _ {n} u q _ {m}
\]

对 N 中的 n 和 m。）

\begin{enumerate}
\def\labelenumi{\alph{enumi})}
\setcounter{enumi}{2}
\item
  用反证法推出 u 是紧的（"皮特定理"）。
\item
  证明 \(\ell_{\mathrm{K}}^{p}(\mathbf{N})\) 到 \(\ell_{\mathrm{K}}^{r}(\mathbf{N})\) 的典范单射不是紧的。
\item
  设 p 和 r 是满足 \(1 \leqslant p \leqslant r \leqslant +\infty\) 的实数。证明 \(\ell_{\mathrm{K}}^{r}(\mathbf{N})\) 同构于 \(\ell_{\mathrm{K}}^{p}(\mathbf{N})\) 当且仅当 p = r。
\end{enumerate}

\begin{enumerate}
\def\labelenumi{\arabic{enumi})}
\setcounter{enumi}{10}
\tightlist
\item
  设 E 和 F 是 K 上的巴拿赫空间。称线性映射 \(u \in \mathcal{L}(\mathrm{E};\mathrm{F})\) 为完全连续的，如果对 E 中每个弱收敛序列 \((x_{n})_{n \in \mathbb{N}}\)，序列 \((u(x_{n}))_{n \in \mathbb{N}}\) 在 F 中收敛。
\end{enumerate}

\begin{enumerate}
\def\labelenumi{\alph{enumi})}
\item
  每个紧线性映射都是完全连续的。
\item
  若 E 的对偶空间 E' 是可数生成的，则每个完全连续线性映射 E → F 都是紧的。
\item
  设 \((u_{n})_{n\in\mathbf{N}}\) 是 \(\ell_{\mathrm{K}}^{1}(\mathbf{N})\) 中元素的一个序列。证明 \((u_{n})\) 收敛当且仅当它弱收敛（"舒尔定理"；参见 EVT, IV, p. 54, 习题 15）。
\end{enumerate}

\(d)\) 由此推出 \(\ell_{\mathrm{K}}^{1}(\mathbf{N})\) 到 \(\ell_{\mathrm{K}}^{2}(\mathbf{N})\) 的单射是完全连续的，但不是紧的。

\begin{enumerate}
\def\labelenumi{\alph{enumi})}
\setcounter{enumi}{4}
\item
  对每个巴拿赫空间，每个连续线性映射 \(\mathrm{E} \to \ell_{\mathrm{K}}^{1}(\mathbf{N})\) 都是完全连续的，且每个连续线性映射 \(\ell_{\mathrm{K}}^{1}(\mathbf{N}) \to \mathrm{E}\) 都是完全连续的。
\item
  若 E 是自反的或其对偶是可数生成的，则每个连续线性映射 \(\mathrm{E} \rightarrow \ell_{\mathrm{K}}^{1}(\mathbf{N})\) 都是紧的。
\end{enumerate}

\begin{enumerate}
\def\labelenumi{\arabic{enumi})}
\setcounter{enumi}{11}
\tightlist
\item
  设 E 是一个无限维自反巴拿赫空间。设 \(u \in \mathcal{L}^{c}(E)\) 是一个紧自同态。假设 u 不是有限秩的。设 \(E_{\sigma}\) 是赋以弱拓扑的空间 E。
\end{enumerate}

证明线性映射 \(^{t}u\) 是 \((\mathrm{E}_{\sigma})'\) 的一个紧自同态。

\begin{enumerate}
\def\labelenumi{\arabic{enumi})}
\setcounter{enumi}{12}
\tightlist
\item
  设 E 是 K 上的一个可数生成的巴拿赫空间。回顾（EVT, IV, p. 70, 习题 14）：E 中元素的一个序列 \((x_{n})_{n\in \mathbf{N}}\) 称为 E 的绍德尔基，如果对每个 \(x\in \mathrm{E}\) ，存在唯一的纯量序列 \((\lambda_n(x))_{n\in \mathbf{N}}\)，使得
\end{enumerate}

\[
x = \sum_ {n = 1} ^ {+ \infty} \lambda_ {n} (x) x _ {n}
\]

其中级数在 E 中收敛。

\begin{enumerate}
\def\labelenumi{\alph{enumi})}
\item
  每个可分希尔伯特空间都具有巴拿赫基；每个空间 \(\ell^{p}(\mathrm{I})\)（其中 I 可数）都具有巴拿赫基。
\item
  由复平面中半径为 1 的开圆盘 D 上有界全纯函数组成的巴拿赫空间 E，赋以范数 \(\| f \| = \sup_{z \in \mathrm{D}} |f(z)|\)，具有绍德尔基。
\item
  每个具有巴拿赫基的巴拿赫空间 E 都是一个逼近空间。更精确地说，存在实数 \(c \geqslant 0\)，使得 \(1_{E}\) 属于 \(\mathcal{L}_{\mathrm{pc}}(\mathrm{E})\) 中由 E 中范数 \(\leqslant c\) 的有限秩自同态组成的集合的闭包。
\end{enumerate}

\begin{enumerate}
\def\labelenumi{\arabic{enumi})}
\setcounter{enumi}{13}
\tightlist
\item
  \begin{enumerate}
  \def\labelenumii{\alph{enumii})}
  \tightlist
  \item
    设 E 和 F 是希尔伯特空间。从 E 到 F 的连续线性映射 u 是紧的，当且仅当 u 的像不含无限维的闭子空间。
  \end{enumerate}
\end{enumerate}

\begin{enumerate}
\def\labelenumi{\alph{enumi})}
\setcounter{enumi}{1}
\tightlist
\item
  a) 中关于从 E 到 F 的紧自同态的刻画并非对每一对巴拿赫空间 (E, F) 都成立。（可利用 p. 121 的习题 5。）
\end{enumerate}

\begin{enumerate}
\def\labelenumi{\arabic{enumi})}
\setcounter{enumi}{14}
\tightlist
\item
  设 \(E_{0}\) 是由函数 \(f: N^{*} \times N^{*} \to C\) 组成的复向量空间。对 \(k \in N^{*}\) ，设 \(\alpha_{k} \in E_{0}\) 是满足下式的函数：
\end{enumerate}

\[
\alpha_ {k} (n, m) = \left\{ \begin{array}{l l} m ^ {k} & \text { if } 1 \leqslant n \leqslant k - 1 \\ n ^ {k} & \text { if } n \geqslant k. \end{array} \right.
\]

设 E 是 \(E_{0}\) 的由满足 \(p_{k}(f)\) 对所有 \(k \in N^{*}\) 有限的函数 f 组成的向量子空间，其中 \(p_{k}\) 由

\[
p _ {k} (f) = \sum_ {n, m} \alpha_ {k} (n, m) | f (n, m) |.
\]

\begin{enumerate}
\def\labelenumi{\alph{enumi})}
\tightlist
\item
  赋以由范数 \((p_{k})_{k\in\mathbf{N}^{*}}\) 定义的拓扑的空间 E 是一个弗雷歇空间，也是一个蒙泰尔空间。它的对偶 E' 等同于 \(E_{0}\) 中由满足下式的函数 g 组成的子空间：对至少一个 \(k\geqslant1\) ，有
\end{enumerate}

\[
\sup _ {m, n} \alpha_ {k} (m, n) ^ {- 1} | g (m, n) | <   + \infty
\]

（参见 EVT, IV, p. 63, 习题 8）。

\begin{enumerate}
\def\labelenumi{\alph{enumi})}
\setcounter{enumi}{1}
\tightlist
\item
  考虑由下式定义的线性映射 \(u\)，它把 E 映到 \(\ell^1 (\mathbf{N}^*)\)：
\end{enumerate}

\[
u (f) = \left(\sum_ {n} f (n, m)\right) _ {m \in \mathbf {N} ^ {*}}.
\]

线性映射 u 是连续的。

\begin{enumerate}
\def\labelenumi{\alph{enumi})}
\setcounter{enumi}{2}
\item
  把 \(\ell^{\infty}(\mathbf{N}^{*})\) 等同于 \(\ell^{1}(\mathbf{N}^{*})\) 的对偶。\(u\) 的转置是映射 \(v\colon \ell^{\infty}(\mathbf{N}^{*})\to \mathrm{E}^{\prime}\subset \mathrm{E}_{0}\)，由 \((x_{n})\mapsto f\) 确定，其中 \(f(n,m) = x_{m}\) 对所有 \((n,m)\in \mathbf{N}^{*}\times \mathbf{N}^{*}\) 成立。
\item
  映射 u 诱导一个同构 \(\mathrm{E}/\mathrm{Ker}(u)\to\ell^{1}(\mathbf{N}^{*})\)（参见出处同上）。
\item
  映射 u 不是紧的。
\item
  u 的转置 v 对于在 \(\ell^{\infty}(\mathbf{N}^{*})\) 的有界子集上收敛的拓扑是紧的。
\item
  映射 v 是分离局部凸空间之间一个紧线性映射的例子，其转置不是紧的。
\end{enumerate}

\begin{enumerate}
\def\labelenumi{\arabic{enumi})}
\setcounter{enumi}{15}
\item
  设 E 是局部凸空间，u 是 E 的一个紧自同态。用 \(E_{b}^{\prime}\) 表示赋有有界收敛拓扑的 E 的对偶空间。线性映射 \(^{t}u \circ ^{t}u\) 是 \(E_{b}^{\prime}\) 的一个紧自同态。
\item
  设 E 是可数型的巴拿赫空间。则 \(\mathcal{L}^{\mathrm{c}}(\mathrm{E})\) 赋以由 \(\mathcal{L}(\mathrm{E})\) 的范数诱导的拓扑后是可数型的，当且仅当对偶空间 \(\mathrm{E}'\) 是可数型的。
\item
  设 E、F、G 是巴拿赫空间，u 是 \(\mathcal{L}^{c}(\mathrm{E};\mathrm{F})\) 的一个元素。
\end{enumerate}

\begin{enumerate}
\def\labelenumi{\alph{enumi})}
\tightlist
\item
  设 v 是 \(\mathcal{L}(\mathrm{E};\mathrm{G})\) 的一个元素。假设对每个 \(\varepsilon > 0\) ，存在 C \textgreater{} 0，使得
\end{enumerate}

\[
\| v (x) \| _ {\mathrm{G}} \leqslant \varepsilon \| x \| _ {\mathrm{E}} + \mathrm{C} \| u (x) \| _ {\mathrm{F}}
\]

对所有 \(x\in \mathrm{E}\) 成立。则 \(v\in \mathcal{L}^{\mathrm{c}}(\mathrm{E};\mathrm{G})\) 。

\begin{enumerate}
\def\labelenumi{\alph{enumi})}
\setcounter{enumi}{1}
\tightlist
\item
  设 w 是 \(\mathcal{L}(\mathrm{F};\mathrm{G})\) 的一个单射元素。则对每个 \(\varepsilon > 0\) ，存在 C \textgreater{} 0，使得
\end{enumerate}

\[
\| u (x) \| _ {\mathrm{F}} \leqslant \varepsilon \| x \| _ {\mathrm{E}} + \mathrm{C} \| (w \circ u) (x) \| _ {\mathrm{G}}
\]

对所有 \(x \in E\) 成立。

\begin{enumerate}
\def\labelenumi{\arabic{enumi})}
\setcounter{enumi}{18}
\tightlist
\item
  设 X 是一个局部紧、可度量化且在无穷远处可数的拓扑空间，\(\mu\) 是 X 上的一个正测度。我们用 A 表示 \(\mu\) 的原子集。设 \(f \in \mathcal{L}^{\infty}(X, \mu)\) 且 \(p \in [1, +\infty]\) 。
\end{enumerate}

证明乘以 \(f\) 在 \(\mathrm{L}^{p}(\mathrm{X},\mu)\) 中定义一个紧自同态，当且仅当下列两个条件被满足：

\begin{enumerate}
\def\labelenumi{(\roman{enumi})}
\item
  集合 \(\{x\in X-A\mid f(x)\neq0\}\) 是 \(\mu\) -可忽略的；
\item
  对每个 \(\varepsilon > 0\) ，存在 A 的有限子集 F，使得 \(|f(a)| \leqslant \varepsilon\) 对所有 \(a \in \mathbf{A} - \mathbf{F}\) 成立。
\end{enumerate}

（考虑乘以 \(f\) 的运算在空间 \(\mathrm{L}^{p}(\mathrm{X}_{\varepsilon},\mu)\) 上的作用，其中 \(\mathrm{X}_{\varepsilon}\) 是由 \(x\in\mathrm{X}\) 且满足 \(|f(x)|\geqslant\varepsilon\) 的元素组成的集合。）

\begin{enumerate}
\def\labelenumi{\arabic{enumi})}
\setcounter{enumi}{19}
\tightlist
\item
  设 X、Y 是紧度量空间，T 是从 \(\mathcal{C}(\mathrm{Y})\) 到 \(\mathcal{C}(\mathrm{X})\) 的一个紧算子。
\end{enumerate}

证明存在一个正测度 \(\mu\)（在 \(X\) 上）以及一个可测映射 \(N\)（从 \(X \times Y\) 到 \(\mathbf{C}\)），使得

\begin{enumerate}
\def\labelenumi{\alph{enumi})}
\item
  对每个 \(y \in Y\) ，映射 \(\mathrm{N}_{y} \colon x \mapsto \mathrm{N}(x, y)\) 是从 \(X\) 到 \(\mathbf{C}\) 的 \(\mu\) -可积映射；
\item
  映射 \(y \mapsto N_{y}\) 是从 Y 到 \(\mathcal{L}^{1}(X, \mu)\) 的连续映射；
\item
  对每个 \(f \in \mathcal{C}(\mathrm{X})\) 和每个 \(y \in \mathrm{Y}\) ，我们有
\end{enumerate}

\[
\mathrm{T} f (y) = \int_ {\mathrm{X}} \mathrm{N} (x, y) f (x) d \mu (x).
\]

\begin{enumerate}
\def\labelenumi{\arabic{enumi})}
\setcounter{enumi}{20}
\tightlist
\item
  我们用 D 表示 C 中以 0 为心、以 1 为半径的开圆盘，用 H \(^{2}\) (D) 表示 D 上的哈代空间（EVT, V, p. 7, 例 5）。我们记 T = R/2πZ，并赋以规范化的哈尔测度。设 j 是由下式定义的从 H \(^{2}\) (D) 到 L \(^{2}\) (T) 的等距：
\end{enumerate}

\[
j (f) (\theta) = \sum_ {n \geqslant 0} \frac {f ^ {(n)} (0)}{n !} e ^ {i n \theta}.
\]

H \(^{2}\) (D) 在 \(j\) 下的像将记为 H \(^{2}\) (T)。

设 \(\widetilde{\mathrm{P}}\) 是 \(\mathrm{L}^2 (\mathbf{T})\) 的以 \(\mathrm{H}^2 (\mathbf{T})\) 为像的正交投影，并设 P 是线性映射 \(f\mapsto \widetilde{\mathrm{P}} (f)\)，它从 \(\mathrm{L}^2 (\mathbf{T})\) 到 \(\mathrm{H}^2 (\mathbf{T})\)。若 \(f\) 是 \(\mathrm{L}^{\infty}(\mathbf{T})\) 的一个元素，设 \(\mathrm{M}_f\) 是自同态 \(g\mapsto fg\)，它作用于 \(\mathrm{L}^2 (\mathbf{T})\)。

证明若 \(f\) 是连续函数，则换位子 \(\widetilde{\mathrm{P}}\mathrm{M}_f - \mathrm{M}_f\widetilde{\mathrm{P}}\) 是一个紧线性映射。（考虑 \(f\) 是三角多项式的特殊情形。）

\begin{enumerate}
\def\labelenumi{\arabic{enumi})}
\setcounter{enumi}{21}
\tightlist
\item
  我们保留上一习题的记号。设 R 是从 \(\mathrm{H}^{2}(\mathbf{T})\) 到 \(\mathrm{L}^{2}(\mathbf{T})\) 的由 \(\mathrm{R}f(\theta)=f(-\theta)\) 定义的连续线性映射。对 \(g\in\mathrm{L}^{\infty}(\mathbf{T})\)，我们令 \(\Gamma_{g}=\mathrm{PM}_{g}\mathrm{R}\) 。
\end{enumerate}

\begin{enumerate}
\def\labelenumi{\alph{enumi})}
\item
  证明线性映射 \(\Gamma\colon g\mapsto\Gamma_{g}\) 是从 \(\mathrm{L}^{\infty}(\mathbf{T})\) 到 \(\mathcal{L}(\mathrm{H}^{2}(\mathbf{T}))\) 的连续映射，范数 \(\leqslant 1\)，且其核 F 等于 \(\mathrm{H}^{2}(\mathbf{T})^{\circ}\cap\mathrm{L}^{\infty}(\mathbf{T})\) 。
\item
  若 \(g\) 是连续函数，则 \(\Gamma_g\) 是 \(\mathrm{H}^2 (\mathbf{T})\) 的一个紧自同态。
\item
  设 \(g \in \mathrm{L}^{\infty}(\mathbf{T})\) 。我们用 \(\alpha_{g}\) 表示 \(\mathrm{H}^{2}(\mathbf{T})\) 上由下式定义的线性形式
\end{enumerate}

\[
\alpha_ {g} (f) = \int_ {\mathbf {T}} \Gamma_ {g} (f) (\theta) d \theta .
\]

设 \(\mathcal{A}\) 是在 \(j\) 下 \(\mathrm{H}^2 (\mathrm{D})\) 中由可解析延拓到 \(\overline{\mathrm{D}}\) 的一个邻域的全纯函数组成的子空间的像。证明若 \(f_{1}\) 和 \(f_{2}\) 属于 \(\mathcal{A}\) ，则

\[
\alpha_ {g} (f _ {1} f _ {2}) = \int_ {\mathbf {T}} \Gamma_ {g} (f _ {1}) (\theta) f _ {2} (- \theta) d \theta .
\]

\(d)\) 设 \(\varphi\) 是 \(\overline{\mathrm{D}}\) 的一个邻域上的全纯函数，在 \(\partial\mathrm{D}\) 上不为零，并设 \(f = j(\varphi|\mathrm{D})\) 。则存在 \(f_{1}\) 和 \(f_{2}\)（都在 \(\mathscr{A}\) 中），使得 \(f = f_{1}f_{2}\)，且

\[
\| f \| _ {\mathrm{L} ^ {1} (\mathbf {T})} = \| f _ {1} \| _ {\mathrm{L} ^ {2} (\mathbf {T})} \| f _ {2} \| _ {\mathrm{L} ^ {2} (\mathbf {T})}.
\]

（设 \(z_{1},\ldots,z_{n}\) 是 \(\varphi\) 在 D 中的零点；则存在一个在 D 的邻域内全纯的函数 \(\psi\)，使得

\[
\varphi (z) = \prod_ {k = 1} ^ {n} \frac {z - z _ {k}}{1 - \overline {{z}} _ {k} z} \psi (z) ^ {2}
\]

对所有 z 成立。）

\begin{enumerate}
\def\labelenumi{\alph{enumi})}
\setcounter{enumi}{4}
\tightlist
\item
  证明对所有 \(f \in \mathrm{H}^{2}(\mathbf{T})\) ，我们有不等式
\end{enumerate}

\[
\left| \alpha_ {g} (f) \right| \leqslant \left\| \Gamma_ {g} \right\| _ {\mathscr {L} \left(\mathrm{H} ^ {2} (\mathbf {T})\right)} \| f \| _ {\mathrm{L} ^ {1} (\mathbf {T})}.
\]

（若 \(\varphi \in \mathrm{H}^2 (\mathrm{D})\) ，则存在 ]0,1[ 中的一个序列 \((r_n)_{n\in \mathbf{N}}\)，使得对每个 \(n\in \mathbf{N}\) ，函数 \(\varphi_{n}\)（由 \(\varphi_{n}(z) = \varphi (r_{n}z)\) 定义）在 \(\partial \mathrm{D}\) 上没有零点，且 \(\lim \varphi_n = \varphi\) 在 \(\mathrm{H}^2 (\mathrm{D}).\) 中

由此推出存在 \(h \in \mathrm{L}^{\infty}(\mathbf{T})\)，使得

\[
\alpha_ {g} (f) = \int_ {\mathbf {T}} f (\theta) h (\theta) d \theta
\]

对所有 \(f \in \mathrm{H}^2(\mathbf{T})\) 成立，且 \(\|h\|_{\mathrm{L}^\infty(\mathbf{T})} \leqslant \| \Gamma_g \|_{\mathcal{L}(\mathrm{H}^2(\mathbf{T}))}\) 。

\(f)\) 证明存在 \(k \in \mathrm{L}^{\infty}(\mathbf{T})\)，使得 \(\Gamma_g = \Gamma_k\) 且 \(\| k \|_{\mathrm{L}^{\infty}(\mathbf{T})} = \| \Gamma_g \|_{\mathcal{L}(\mathrm{H}^2(\mathbf{T}))}\) 。

于是映射 \(\widehat{\Gamma}\colon\mathrm{L}^{\infty}(\mathbf{T})/\mathrm{F}\to\mathcal{L}(\mathrm{H}^{2}(\mathbf{T}))\) 是一个等距（“内哈里定理”）。

\(g)\) 设 \(g\) 是 \(\mathbf{T}\) 上的连续函数。则我们有 \(d(g, \mathrm{F}) = d(g, \mathrm{F} \cap \mathcal{C}(\mathbf{T}))\) 。（对 \(\mathrm{R} > 1\) ，用下式定义线性映射 \(\mathrm{P}_{\mathrm{R}}\)，即从 \(\mathrm{L}^{\infty}(\mathbf{T})\) 到 \(\mathcal{C}(\mathbf{T})\) 的映射：

\[
\mathrm{P} _ {\mathrm{R}} f (\theta) = \int_ {\mathbf {T}} \frac {\mathrm{R} ^ {2} - 1}{| \mathrm{R} e ^ {i (\theta - \varphi)} - 1 | ^ {2}} f (\varphi) d \varphi .
\]

则有 \(\| \mathrm{P}_{\mathrm{R}}f\|_{\mathcal{C}(\mathbf{T})}\leqslant \| f\|_{\mathrm{L}^{\infty}(\mathbf{T})}\)，且若 \(f\in \mathcal{C}(\mathbf{T})\) ，则 \(\mathrm{P_R}f\to f\) 在 \(\mathcal{C}(\mathbf{T})\) 中成立（当 \(\mathrm{R}\rightarrow 1.\)）。

由此推出 \(\Gamma(\mathcal{C}(\mathbf{T}))\) 是 \(\mathcal{L}(\mathrm{H}^{2}(\mathbf{T}))\) 的一个闭子空间（“萨拉森定理”）。

\begin{enumerate}
\def\labelenumi{\alph{enumi})}
\setcounter{enumi}{7}
\tightlist
\item
  设 \(g \in \mathrm{L}^{\infty}(\mathbf{T})\) 满足 \(\Gamma_{g}\) 是一个紧自同态。对 \(m \in Z\)，我们令 \(e_{m}(\theta) = e^{im\theta}\) 。证明 \(\|\Gamma_{ge_{-m}}\|_{\mathcal{L}(\mathrm{H}^{2}(\mathbf{T}))}\) 当 m 趋于 \(+\infty\) 时趋于零，然后证明存在 F 中的一个序列 \((h_{m})_{m \in \mathbf{N}}\)，使得序列 \((e_{m}h_{m})_{m \in \mathbf{N}}\) 在 \(\mathrm{L}^{\infty}(\mathbf{T})\) 中收敛于 g。由此推出在 g 模 F 的类中存在一个连续函数（“哈特曼定理”）。
\end{enumerate}

\begin{enumerate}
\def\labelenumi{\arabic{enumi})}
\setcounter{enumi}{22}
\item
  设 B 是一个无限维巴拿赫空间。我们用 E 表示 B 的强对偶，用 F 表示赋有拓扑 \(\sigma(\mathrm{B}^{\prime},\mathrm{B})\) 的 B 的对偶。证明 B 的对偶的恒同映射从 E 到 F 是紧的，但其转置从 \(\mathrm{F}_{b}^{\prime}\) 到 \(\mathrm{E}_{b}^{\prime}\) 不是紧的。
\item
  设 E 是一个实或复的巴拿赫空间，\(u \in \mathcal{L}(\mathrm{E})\) 。假设序列 \((u^{n})_{n \in \mathbf{N}}\) 有界，且 u 从 E 到赋有弱拓扑 \(\sigma(\mathrm{E}, \mathrm{E}')\) 的 E 是紧的。对每个整数 \(N \geqslant 1\)，我们令
\end{enumerate}

\[
v _ {\mathrm{N}} = \frac {1}{\mathrm{N}} (1 _ {\mathrm{E}} + u + \dots + u ^ {\mathrm{N-1}}) \in \mathcal {L} (\mathrm{E}).
\]

\begin{enumerate}
\def\labelenumi{\alph{enumi})}
\item
  设 F 是 \(1_{\mathrm{E}} - u\) 的像的闭包。对每个 \(x \in \mathrm{F}\) ，我们有 \(v_{\mathrm{N}}(x) \to 0\)（在 E 中，当 \(\mathrm{N} \to +\infty\) 时）。
\item
  设 \(x \in E\) 。序列 \((v_{\mathrm{N}}(x))\) 收敛到一个元素 \(y \in E\)，使得 \(u(y) = y\) 。（存在 \((v_{\mathrm{N}}(x))\) 的一个子序列，它弱收敛到 \(y \in E\) ；利用哈恩–巴拿赫定理证明 \(x - y \in F\) 。）
\item
  映射 \(w \colon \mathrm{E} \to \mathrm{E}\) 满足 \(w(x)\) 是上一问题给出的元素 \(y\)，它是一个连续投影，其像是 \(u\) 的对应于特征值 1 的特征空间。
\end{enumerate}

\(d)\) 设 \(z \in \mathbf{C}\) 满足 \(|z| = 1\) 。由下式定义的 \(\mathcal{L}(\mathrm{E})\) 中的序列

\[
v _ {z, \mathrm{N}} = \frac {1}{\mathrm{N}} (1 _ {\mathrm{E}} + z ^ {- 1} u \dots + z ^ {- (\mathrm{N-1})} u ^ {\mathrm{N-1}})
\]

在赋有逐点收敛拓扑的 \(\mathcal{L}(\mathrm{E})\) 中收敛。用 \(w_{z}\) 表示它的极限。则 \(w_{z}\) 是一个投影，我们有 \(w_{z}w = ww_{z} = zw_{z}\)，以及 \(w_{z_1}w_{z_2} = 0\)（若 \(z_{1} \neq z_{2}\)）。

\begin{enumerate}
\def\labelenumi{\arabic{enumi})}
\setcounter{enumi}{24}
\tightlist
\item
  \begin{enumerate}
  \def\labelenumii{\alph{enumii})}
  \tightlist
  \item
    设 \(N \geqslant 1\) 是一个整数。我们用 \(G_{N}\)（相应地，\(H_{N}\)）表示集合 \(\{-1,1\}^{N}\)（相应地，集合 \(\{-1,2\}^{N}\)）。我们用 \(\mu_{N}\)（相应地，\(\nu_{N}\)）表示 \(G_{N}\) 上（相应地，\(H_{N}\) 上）总质量为 1 的正测度，使得对每个 \((t_{i})_{1 \leqslant i \leqslant N}\)（属于 \(G_{N}\)，相应地，属于 \(H_{N}\)），有 \(\mu_{\mathrm{N}}((t_{i})_{1 \leqslant i \leqslant N}) = \frac{1}{2^{\mathrm{N}}}\) ，相应地，
  \end{enumerate}
\end{enumerate}

\[
\nu_ {\mathrm{N}} ((t _ {i}) _ {1 \leqslant i \leqslant \mathrm{N}}) = \left(\frac {1}{3}\right) ^ {j} \left(\frac {2}{3}\right) ^ {\mathrm{N} - j},
\]

其中 j 是使得 \(t_{i}=2\) 的指标 i 的个数。

存在一个实数 \(c > 0\)，使得对每个由复数组成的族 \((a_{i})_{1\leqslant i\leqslant N}\)，我们有

\[
\mu_ {\mathrm{N}} \Big (\Big \{(t _ {i}) _ {1 \leqslant i \leqslant \mathrm{N}} \in \mathrm{G} _ {\mathrm{N}} | \Big | \sum_ {i = 1} ^ {\mathrm{N}} a _ {i} t _ {i} \Big | > c (\log \mathrm{N}) ^ {1 / 2} \Big (\sum_ {i = 1} ^ {\mathrm{N}} | a _ {i} | ^ {2} \Big) ^ {1 / 2} \Big \} \Big) <   \frac {c}{\mathrm{N} ^ {3}},
\]

以及

\[
\nu_ {\mathrm{N}} \Big (\Big \{(t _ {i}) _ {1 \leqslant i \leqslant \mathrm{N}} \in \mathrm{H} _ {\mathrm{N}} | \left| \sum_ {i = 1} ^ {\mathrm{N}} a _ {i} t _ {i} \right| > c (\log \mathrm{N}) ^ {1 / 2} \Big (\sum_ {i = 1} ^ {\mathrm{N}} | a _ {i} | ^ {2} \Big) ^ {1 / 2} \Big \} \Big) <   \frac {c}{\mathrm{N} ^ {3}}.
\]

（我们可以假设 \(a_{i} \in \mathbf{R}\) 对所有 \(i\) 成立，且 \(\sum |a_{i}|^{2} = 1\) 。对第一个不等式，证明并利用上界估计

\[
\int_ {G _ {N}} \exp \left(\lambda \left| \sum_ {i = 1} ^ {N} a _ {i} t _ {i} \right|\right) d \mu_ {N} \leqslant 2 \exp (\lambda^ {2})
\]

对每个实数 \(\lambda > 0\) 成立；对第二个上界估计，利用不等式 \(\frac{1}{3}(e^{2x} + 2e^{-x}) \leqslant e^{2x^2}\)（对所有 \(x \in \mathbf{R}\)）以类似的方式推理。）

\begin{enumerate}
\def\labelenumi{\alph{enumi})}
\setcounter{enumi}{1}
\tightlist
\item
  对 \(k \in N\)，令 \(\mathrm{A}_{k} = \mathbf{Z}/(3 \cdot 2^{k})\mathbf{Z}\) 。存在一个不依赖于 k 的实数 \(c' > 0\)，以及一个函数 \(f_{k} \colon \widehat{\mathrm{A}}_{k} \to \{-1, 2\}\)，使得
\end{enumerate}

\[
\sum_ {\chi \in \widehat {\mathrm{A}} _ {k}} f _ {k} (\chi) = 0
\]

且对所有 \(x \in A_{k}\) ，有

\[
\Big | \sum_ {\chi \in \widehat {\mathrm{A}} _ {k}} \chi (x) f _ {k} (\chi) \Big | \leqslant c ^ {\prime} (k + 1) ^ {1 / 2} 2 ^ {k / 2}.
\]

（利用 \(a\)），确定满足第二个条件的函数 \(f_{k}\) 的存在性，然后修改 \(f_{k}\) 以得到第一个条件。）

我们用 \(B_{k}\)（相应地，\(C_{k}\)）表示特征标 \(\chi\)（\(A_{k}\) 上的）中满足 \(f_{k}(\chi)=2\)（相应地，\(f_{k}(\chi)=-1\)）者所成的集合。我们有 \(\mathrm{Card(B_{k})}=2^{k}\) 和 \(\mathrm{Card(C_{k})}=2^{k+1}\) 。

对每个 \(k \geqslant 1\) ，我们固定一个双射 \(\iota_k \colon B_k \to C_{k-1}\) 。

我们用 A 表示诸 \(A_{k}\)（\(k \in N\)）的不交并；我们赋给 A 离散拓扑，并记 \(\mathrm{E} = \mathcal{C}_{b}(\mathrm{A})\) 。

我们固定一个函数族 \(\varepsilon = (\varepsilon_k)_{k\geqslant \mathbf{N}}\)，其中的函数 \(\varepsilon_{k}\colon \mathrm{B}_{k}\to \{-1,1\}\)。然后我们定义函数 \(e_{k,\chi}\in \mathrm{E}\)（对 \(k\in \mathbf{N}\) 和 \(\chi \in \mathrm{B}_k\)）如下：

\[
e _ {k, \chi} (x) = \left\{ \begin{array}{l l} \iota_ {k} (\chi) (x) & \text { if } k \geqslant 1 \text {   and   } x \in \mathrm{A} _ {k - 1}, \\ \varepsilon_ {k} (\chi) \chi (x) & \text { if } x \in \mathrm{A} _ {k} \\ 0 & \text { otherwise. } \end{array} \right.
\]

我们用 \(E_{\varepsilon}\) 表示由函数 \((e_{k,\chi})_{k\in\mathbf{N},\chi\in\mathbf{B}_{k}}\) 生成的 E 的闭子空间。这是一个可数型的复巴拿赫空间。习题的其余部分旨在证明：对 \(\varepsilon\) 的适当选取，空间 \(\mathrm{E}_{\varepsilon}\) 不具有逼近性质；它不拥有巴拿赫基。

\begin{enumerate}
\def\labelenumi{\alph{enumi})}
\setcounter{enumi}{2}
\tightlist
\item
  对 \(k \in N\) 和 \(\chi \in B_{k}\)，存在一个连续线性形式 \(\lambda_{k,\chi} \in E_{\varepsilon}'\)，使得对每个 \(f \in E_{\varepsilon}\)，
\end{enumerate}

\[
\lambda_ {k, \chi} (f) = \frac {1}{\operatorname{Card} (\mathrm{A} _ {k})} \sum_ {x \in \mathrm{A} _ {k}} \varepsilon_ {k} (\chi) \chi (x ^ {- 1}) f (x).
\]

对 \(k \geqslant 1\) 和 \(f \in \mathrm{E}_{\varepsilon}\) ，我们有

\[
\lambda_ {k, \chi} (f) = \frac {2}{\operatorname{Card} (\mathrm{A} _ {k})} \sum_ {x \in \mathrm{A} _ {k}} \iota_ {k} (\chi) (x ^ {- 1}) f (x).
\]

\begin{enumerate}
\def\labelenumi{\alph{enumi})}
\setcounter{enumi}{3}
\tightlist
\item
  对每个 \(k \in N\) ，用 \(\beta_{k}\) 表示 \(\mathcal{L}(E_{\varepsilon})\) 上由下式定义的连续线性形式
\end{enumerate}

\[
\beta_ {k} (u) = \frac {1}{2 ^ {k}} \sum_ {\chi \in \mathrm{B} _ {k}} \lambda_ {k, \chi} (u (e _ {k, \chi})).
\]

我们有 \(\beta_{k}(1_{\mathrm{E}_{\varepsilon}}) = 1\)

\begin{enumerate}
\def\labelenumi{\alph{enumi})}
\setcounter{enumi}{4}
\tightlist
\item
  设 \(u \in \mathcal{L}(\mathrm{E}_{\varepsilon})\) 且 \(k \in \mathbf{N}\) 。我们有
\end{enumerate}

\[
\beta_ {k + 1} (u) - \beta_ {k} (u) = \frac {1}{\operatorname{Card} \left(\mathrm{A} _ {k}\right)} \sum_ {x \in \mathrm{A} _ {k}} u \left(\varphi_ {k, x}\right) (x)
\]

其中，对 \(x \in \mathrm{A}_k\) ，我们令

\[
\varphi_ {k, x} = \frac {1}{2 ^ {k + 1}} \sum_ {\chi \in B _ {k + 1}} \iota_ {k + 1} (\chi) (x ^ {- 1}) e _ {k + 1, \chi} - \frac {1}{2 ^ {k}} \sum_ {\chi \in B _ {k}} \varepsilon_ {k} (\chi) \chi (x ^ {- 1}) e _ {k, \chi}.
\]

我们有 \(|\varphi_{k,x}(y)| \leqslant c'(k + 1)^{1/2}2^{-k/2}\)，对 \((x,y) \in \mathrm{A}_k^2\) 成立。

\begin{enumerate}
\def\labelenumi{\alph{enumi})}
\setcounter{enumi}{5}
\item
  存在一个实数 \(c'' \geqslant c'\)，使得对每个 \(k \geqslant 1\) ，存在一个函数 \(\varepsilon_k \colon B_k \to \{-1,1\}\)，满足 \(|\varphi_{k,x}(y)| \leqslant c''3(k+1)^{1/2}2^{-k/2}\) 对 \((x,y) \in A_k \times A_{k-1}\) 成立。（应用 a) 的第一个不等式。）我们现在假设 \(\varepsilon = (\varepsilon_k)\) 满足这些性质。
\item
  对每个 \(k \in \mathbf{N}\) 和 \(x \in \mathrm{A}_k\) ，我们有 \(\| \varphi_{k,x} \| \leqslant c''(k + 1)^{1/2} 2^{-k/2}\) 。
\item
  设 K 是 \(E_{\varepsilon}\) 的子集，其元素为 \(e_{0,\chi_{0}}\)（其中 \(\chi_{0}\) 是 \(B_{0}\) 的唯一元素）以及函数 \((k+1)^{2}\varphi_{k,x}\)（\(k\in N\)，\(x\in A_{k}\)）。集合 K 在 \(E_{\varepsilon}\) 中是相对紧的，且我们有
\end{enumerate}

\[
| \beta_ {k + 1} (u) - \beta_ {k} (u) | \leqslant \frac {1}{(k + 1) ^ {2}} \sup _ {x \in K} \| u (x) \|
\]

对所有 \(k \in \mathbf{N}\) 和每个 \(u \in \mathcal{L}(\mathrm{E}_{\varepsilon})\) 成立。

\begin{enumerate}
\def\labelenumi{\roman{enumi})}
\tightlist
\item
  序列 \((\beta_{k})_{k\in \mathbf{N}}\) 在赋有逐点收敛拓扑的空间 \(\mathcal{L}(\mathrm{E}_{\varepsilon})\) 的对偶中收敛。用 \(\beta\) 表示它的极限。对每个 \(u\in \mathcal{L}(\mathrm{E}_{\varepsilon})\) ，我们有 \(|\beta (u)|\leqslant 3\sup_{x\in \mathrm{K}}\| u(x)\|\) 。
\end{enumerate}

\begin{enumerate}
\def\labelenumi{\alph{enumi})}
\setcounter{enumi}{9}
\tightlist
\item
  我们有 \(\beta(1_{\mathrm{E}_{\varepsilon}})=1\)，且 \(\beta(u)=0\)（若 \(u\) 是有限秩的）。
\end{enumerate}

\(k)\) 空间 \(\mathrm{E}_{\varepsilon}\) 不具有逼近性质。

\begin{enumerate}
\def\labelenumi{\alph{enumi})}
\setcounter{enumi}{11}
\tightlist
\item
  设 \(p \in]2, +\infty[\) 。我们赋给 A 计数测度，并用 \(E_{\varepsilon}^{p}\) 表示 \(L^{p}(A)\) 的由 \((e_{k,\chi})_{k \in \mathbf{N}, \chi \in \mathrm{B}_{k}}\) 生成的闭子空间。空间 \(E_{\varepsilon}^{p}\) 不具有逼近性质。（验证 \(\varphi_{k,x}\) 的支撑的基数 \(\leqslant 21 \cdot 2^{k}\)，并由此推出 \(\|\varphi_{k,x}\|_{p} = \mathrm{O}((k+1)^{1/2}2^{-k(p-2)/2p})\) 对 \(k \in N\) 和 \(x \in A_{k}\) 成立。）
\end{enumerate}

\begin{enumerate}
\def\labelenumi{\arabic{enumi})}
\setcounter{enumi}{25}
\tightlist
\item
  设 E 和 F 是巴拿赫空间，u 是从 E 到 F 的一个连续线性映射。
\end{enumerate}

\begin{enumerate}
\def\labelenumi{\alph{enumi})}
\tightlist
\item
  证明下列条件等价：
\end{enumerate}

\begin{enumerate}
\def\labelenumi{(\roman{enumi})}
\item
  映射 u 是紧的；
\item
  存在一个序列 \((x_n')_n\)（在 E 的强对偶 \(\mathrm{E}'\) 中趋于 0），满足 \(\| u(x)\| \leqslant \sup |\langle x,x_n' \rangle|\) 对所有 \(x \in \mathrm{E}\) 成立；
\item
  存在巴拿赫空间 \(c_{0}\) 的一个闭子空间 L 以及紧线性映射 \(f: E \to L\) 和 \(g: L \to F\)，使得 \(u = g \circ f\)（若 u 是紧的，对 F 的单位球的极集在 \(^{t}u\) 下的像应用 EVT, IV, p. 24 的推论 1。在假设 (ii) 下，设 \((\lambda_{n})\) 是 \(c_{0}\) 的一个元素，使得序列 \((\lambda_{n}^{-1}x_{n}^{\prime})\) 趋于 0；由 \(f(x) = (\langle x, \lambda_{n}^{-1}x_{n}^{\prime} \rangle)_{n}\) 定义 f，并取 L 为 f 的像的闭包。）
\end{enumerate}

\begin{enumerate}
\def\labelenumi{\alph{enumi})}
\setcounter{enumi}{1}
\tightlist
\item
  若 \(u\) 是紧的，我们有 \(\|u\| = \inf(\sup_n \|x_n'\|)\)，其中下确界取遍满足 (ii) 的序列 \((x_n')\) 所成的集合，且 \(\|u\| = \inf(\|f\|\|g\|)\) ，下确界取遍满足 (iii) 的条件的三元组 (L, \(f,g\) ) 所成的集合。
\end{enumerate}

\begin{enumerate}
\def\labelenumi{\arabic{enumi})}
\setcounter{enumi}{26}
\tightlist
\item
  设 E 和 F 是局部凸空间。映射 \(u \in \mathcal{L}(\mathrm{E};\mathrm{F})\) 称为弱紧的，如果它是从 E 到赋有拓扑 \(\sigma(\mathrm{F},\mathrm{F}')\) 的 F 的紧映射。
\end{enumerate}

\begin{enumerate}
\def\labelenumi{\alph{enumi})}
\item
  弱紧映射构成一个子空间 \(\mathcal{L}^{\mathrm{fc}}(\mathrm{E};\mathrm{F})\)，它是 \(\mathcal{L}(\mathrm{E};\mathrm{F})\) 的子空间。设 \(\mathrm{E}_1\) 、\(\mathrm{F}_1\) 是局部凸空间，\(f\in \mathcal{L}(\mathrm{E}_1;\mathrm{E})\)，\(g\in \mathcal{L}(\mathrm{F};\mathrm{F}_1)\) 且 \(u\in \mathcal{L}(\mathrm{E};\mathrm{F})\)；若 \(u\) 是弱紧的，则 \(g\circ u\circ f\) 亦然。
\item
  我们有 \(\mathcal{L}^{c}(\mathrm{E};\mathrm{F})\subset\mathcal{L}^{\mathrm{fc}}(\mathrm{E};\mathrm{F})\) 。若 E 是无限维自反巴拿赫空间，则 E 的恒同映射是弱紧的但不是紧的。
\end{enumerate}

\begin{enumerate}
\def\labelenumi{\arabic{enumi})}
\setcounter{enumi}{27}
\tightlist
\item
  设 E 和 F 是局部凸空间，且 \(u \in \mathcal{L}(\mathrm{E};\mathrm{F})\) 。
\end{enumerate}

\begin{enumerate}
\def\labelenumi{\alph{enumi})}
\tightlist
\item
  证明下列性质等价：
\end{enumerate}

\begin{enumerate}
\def\labelenumi{(\roman{enumi})}
\item
  E 的每个有界子集在 \(u\) 下的像对于弱拓扑是相对紧的；
\item
  映射 \(^{tt}u\) 把 F'' 映到 E 中。
\end{enumerate}

\begin{enumerate}
\def\labelenumi{\alph{enumi})}
\setcounter{enumi}{1}
\tightlist
\item
  若 F 是拟完备的，证明这些性质等价于：
\end{enumerate}

\begin{enumerate}
\def\labelenumi{(\roman{enumi})}
\setcounter{enumi}{2}
\tightlist
\item
  映射 \(^{t}u\) 把 F' 的每个等度连续子集映到 E' 的一个对于 \(\sigma(\mathrm{E}',\mathrm{E}'')\) 相对紧的子集。
\end{enumerate}

\begin{enumerate}
\def\labelenumi{\alph{enumi})}
\setcounter{enumi}{2}
\tightlist
\item
  若 E 和 F 是巴拿赫空间，则 (i) 至 (iii) 也等价于
\end{enumerate}

(i') 映射 u 是弱紧的；

(iii') 映射 \(^{t}u\colon F_{b}^{\prime}\to E_{b}^{\prime}\) 是弱紧的。

\begin{enumerate}
\def\labelenumi{\arabic{enumi})}
\setcounter{enumi}{28}
\tightlist
\item
  \begin{enumerate}
  \def\labelenumii{\alph{enumii})}
  \tightlist
  \item
    证明巴拿赫空间 \(\ell^{1}(\mathbf{N})\) 的每个弱紧子集是紧的（参见 EVT, IV, p. 69, 习题 11）。特别地，从局部凸空间 E 到 \(\ell^{1}(\mathbf{N})\) 的每个弱紧线性映射是紧的。
  \end{enumerate}
\end{enumerate}

\begin{enumerate}
\def\labelenumi{\alph{enumi})}
\setcounter{enumi}{1}
\tightlist
\item
  设 E 是巴拿赫空间。由 a) 推出，从 \(c_{0}\) 到 E 的每个弱紧线性映射是紧的。
\end{enumerate}

\begin{enumerate}
\def\labelenumi{\arabic{enumi})}
\setcounter{enumi}{29}
\tightlist
\item
  \begin{enumerate}
  \def\labelenumii{\alph{enumii})}
  \tightlist
  \item
    设 E 和 F 是巴拿赫空间，\(u \in \mathcal{L}(\mathrm{E};\mathrm{F})\) 。证明下列性质等价：
  \end{enumerate}
\end{enumerate}

\begin{enumerate}
\def\labelenumi{(\roman{enumi})}
\item
  E 的每个弱紧子集在 \(u\) 下的像是（强）紧的；
\item
  E 的每个弱相对紧子集在 \(u\) 下的像在 F 中相对紧；
\item
  E 中每个弱收敛序列在 \(u\) 下的像在 F 中强收敛。
\end{enumerate}

（利用施穆良定理，EVT, IV, p. 36, 定理 2。）

如果 u 满足这些性质，就说 u 是半紧的 \(^{(3)}\) 。

\begin{enumerate}
\def\labelenumi{\alph{enumi})}
\setcounter{enumi}{1}
\item
  设 E 和 F 是巴拿赫空间。证明由半紧映射组成的集合 \(\mathcal{L}^{\mathrm{sc}}(\mathrm{E};\mathrm{F})\) 是 \(\mathcal{L}(\mathrm{E};\mathrm{F})\) 的一个闭子空间。
\item
  设 E 和 F 是巴拿赫空间，\(u \in \mathcal{L}(\mathrm{E};\mathrm{F})\) 。设 \(E_{1}\) 、\(F_{1}\) 是巴拿赫空间，\(f \in \mathcal{L}(\mathrm{E}_{1};\mathrm{E})\) 且 \(g \in \mathcal{L}(\mathrm{F};\mathrm{F}_{1})\) 。若 \(u: E \to F\) 是半紧的，则 \(g \circ u \circ f\) 亦然。
\item
  设 E 和 F 是巴拿赫空间，\(u \in \mathcal{L}(\mathrm{E};\mathrm{F})\) 。若 u 是紧的，则它是半紧的。反之，若 E 是自反的且 u 是半紧的，则 u 是紧的。
\item
  设 E 和 F 是巴拿赫空间。若 E 的对偶空间 E′ 是可数型的，则每个完全连续线性映射 E → F 都是紧的。
\item
  设 I 是无限集；\(\ell^{1}(\mathbf{I})\) 的恒同映射是半紧的，但不是弱紧的（参见 EVT, IV, p. 54, 习题 15）。g) \(\ell_{\mathrm{K}}^{1}(\mathbf{N})\) 到 \(\ell_{\mathrm{K}}^{2}(\mathbf{N})\) 的单射是半紧的但不是紧的（参见出处同上）。
\item
  设 E 是巴拿赫空间。每个连续线性映射 \(E \rightarrow \ell_{K}^{1}(\mathbf{N})\) 都是半紧的，且每个连续线性映射 \(\ell_{\mathrm{K}}^{1}(\mathbf{N}) \rightarrow \mathrm{E}\) 都是半紧的。
(3) 一些作者称它们为完全连续映射。
\item
  设 E 是巴拿赫空间。若 E 是自反的或其对偶是可数型的，则每个连续线性映射 \(\mathrm{E} \to \ell_{\mathrm{K}}^{1}(\mathbf{N})\) 都是紧的。
\end{enumerate}

\begin{enumerate}
\def\labelenumi{\arabic{enumi})}
\setcounter{enumi}{30}
\tightlist
\item
  设 E 是巴拿赫空间。
\end{enumerate}

\begin{enumerate}
\def\labelenumi{\alph{enumi})}
\tightlist
\item
  证明下列条件等价：
\end{enumerate}

\begin{enumerate}
\def\labelenumi{(\roman{enumi})}
\item
  对任何巴拿赫空间 F，从 E 到 F 的每个弱紧映射都是半紧的。
\item
  若 \((x_{n})\) 是 E 中对于拓扑 \(\sigma(\mathrm{E},\mathrm{E}^{\prime})\) 趋于 0 的序列，而 \((x_{n}^{\prime})\) 是 \(E^{\prime}\) 中对于拓扑 \(\sigma(\mathrm{E}^{\prime},\mathrm{E}^{\prime\prime})\) 趋于 0 的序列，则序列 \(\langle x_{n},x_{n}^{\prime}\rangle\) 趋于 0。
\end{enumerate}

（为证明 (i) 蕴含 (ii)，取 \(F = c_0\) 。若 (i) 不被满足，则存在从 E 到某个巴拿赫空间 F 的一个弱紧映射 \(u\) 以及 E 中弱趋于 0 的一个序列 \((x_n)\)，使得 \(\inf(\|u(x_n)\|) > 0\) ；构造一个序列 \((y_n)\)（在 \(E'\) 的单位球中），它对于 \(\sigma(E', E'')\) 收敛且满足 \(\langle u(x_n), y_n \rangle = \|u(x_n)\|\) ，并取 \(x_n' = y_n - \lim y_n\) 。）

这些条件特别在下列情形成立：\(E = \mathcal{C}(T)\)（其中 T 是紧空间），以及 \(E = L^{1}(X, \mu)\) 或 \(E = L^{\infty}(X, \mu)\)（其中 X 是局部紧空间，\(\mu\) 是 X 上的一个测度）（INT, VI, § 2, 习题 23）。b) 证明具有上述性质的巴拿赫空间不拥有无限维的自反直因子。

\begin{enumerate}
\def\labelenumi{\alph{enumi})}
\setcounter{enumi}{2}
\tightlist
\item
  设 H 是自反巴拿赫空间，B 是 H' 的赋有弱拓扑的单位球。证明 H 可以等同于 \(\mathcal{C}(\mathrm{B})\) 的一个闭的非直接子空间。
\end{enumerate}

\begin{enumerate}
\def\labelenumi{\arabic{enumi})}
\setcounter{enumi}{31}
\item
  一个拥有绍德尔基的可数型巴拿赫空间（EVT, IV, p. 70, 习题 14）具有逼近性质。
\item
  \begin{enumerate}
  \def\labelenumii{\alph{enumii})}
  \tightlist
  \item
    证明逼近空间序列的严格归纳极限是逼近空间。
  \end{enumerate}
\end{enumerate}

\begin{enumerate}
\def\labelenumi{\alph{enumi})}
\setcounter{enumi}{1}
\tightlist
\item
  证明一族逼近空间的局部凸直和是逼近空间。
\end{enumerate}

\begin{enumerate}
\def\labelenumi{\arabic{enumi})}
\setcounter{enumi}{33}
\tightlist
\item
  设 E 是巴拿赫空间。证明下列性质等价：
\end{enumerate}

\begin{enumerate}
\def\labelenumi{(\roman{enumi})}
\item
  空间 E' 是逼近空间；
\item
  对每个巴拿赫空间 F，\(\mathcal{L}^{\mathrm{f}}(\mathrm{E};\mathrm{F})\) 在 \(\mathcal{L}(\mathrm{E};\mathrm{F})\) 中的闭包是 \(\mathcal{L}^{\mathrm{c}}(\mathrm{E};\mathrm{F})\) ；
\item
  对 \(c_{0}\) 的每个闭子空间 F，\(\mathcal{L}^{\mathrm{f}}(\mathrm{E};\mathrm{F})\) 在 \(\mathcal{L}(\mathrm{E};\mathrm{F})\) 中的闭包是 \(\mathcal{L}^{\mathrm{c}}(\mathrm{E};\mathrm{F})\) 。
\end{enumerate}

\begin{enumerate}
\def\labelenumi{\arabic{enumi})}
\setcounter{enumi}{34}
\tightlist
\item
  设 E 是巴拿赫空间。我们称 E 具有度量逼近性质，如果 E 的恒同映射属于范数 \(\leqslant 1\) 的元素（\(\mathcal{L}^{\mathrm{f}}(\mathrm{E})\) 中的）组成的集合在赋有紧收敛拓扑的空间 \(\mathcal{L}(\mathrm{E})\) 中的闭包。
\end{enumerate}

证明下列性质等价：

\begin{enumerate}
\def\labelenumi{(\roman{enumi})}
\item
  空间 E 具有度量逼近性质；
\item
  \(\mathcal{L}(\mathrm{E})\) 的单位球是范数 \(\leqslant 1\) 的元素（\(\mathcal{L}^{\mathrm{f}}(\mathrm{E})\) 中的）组成的集合对于紧收敛拓扑的闭包；
\item
  对每个巴拿赫空间 F，\(\mathcal{L}(\mathrm{E};\mathrm{F})\) 的单位球是范数 \(\leqslant 1\) 的元素（\(\mathcal{L}^{\mathrm{f}}(\mathrm{E};\mathrm{F})\) 中的）组成的集合对于紧收敛拓扑的闭包；
\item
  对每个巴拿赫空间 F，\(\mathcal{L}(\mathrm{F};\mathrm{E})\) 的单位球是范数 \(\leqslant 1\) 的元素（\(\mathcal{L}^{\mathrm{f}}(\mathrm{F};\mathrm{E})\) 中的）组成的集合对于紧收敛拓扑的闭包。
\end{enumerate}

\begin{enumerate}
\def\labelenumi{\arabic{enumi})}
\setcounter{enumi}{35}
\tightlist
\item
  \begin{enumerate}
  \def\labelenumii{\alph{enumii})}
  \tightlist
  \item
    设 X 是局部紧拓扑空间。证明空间 \(\mathcal{C}_{0}(\mathrm{X})\) 具有度量逼近性质（改造 III, p. 21 命题 13 的证明）。
  \end{enumerate}
\end{enumerate}

\begin{enumerate}
\def\labelenumi{\alph{enumi})}
\setcounter{enumi}{1}
\tightlist
\item
  设 X 是赋有正测度的局部紧拓扑空间，且 \(p \in [1, +\infty]\) 。证明空间 \(L^{p}(X)\) 具有度量逼近性质（同样的方法）。
\end{enumerate}

\begin{enumerate}
\def\labelenumi{\arabic{enumi})}
\setcounter{enumi}{36}
\tightlist
\item
  设 X 是 \(C^{r}\) 类（\(1 \leqslant r \leqslant \infty\)）微分流形，仿紧且局部有限维，并设 E 是 X 上的 \(C^{r}\) 类向量丛，局部有限维。证明空间 \(\mathcal{S}^{r}(\mathrm{X};\mathrm{E})\) 具有逼近性质。
\end{enumerate}

\exercisesection{}

\begin{enumerate}
\def\labelenumi{\arabic{enumi})}
\tightlist
\item
  设 \(I = [a, b]\) 是 R 的一个紧区间。对每个连续函数 \(f: I \to C\) 和每个 \(y \in I\)，令
\end{enumerate}

\[
\mathrm{V} (f) (y) = \int_ {a} ^ {y} f (x) d x.
\]

证明从 \(\mathcal{C}(\mathrm{I})\) 到 \(\mathcal{C}(\mathrm{I})\) 的线性映射 V 是紧的（“沃尔泰拉算子”）。

\begin{enumerate}
\def\labelenumi{\arabic{enumi})}
\setcounter{enumi}{1}
\tightlist
\item
  设 X = {]}0, 1{[}，μ 是 X 上的勒贝格测度。用下式定义一个函数序列 \((f_{n})_{n \geqslant 1}\)（X → R）：
\end{enumerate}

\[
f _ {2 n} (x) = \left\{ \begin{array}{l l} \sqrt {2} \sin (2 n \pi x) & \text { if } 0 <   x <   1 / 2 \\ - 2 ^ {n} & \text { if } 1 - 2 ^ {- 2 n} <   x <   1 - 2 ^ {- 2 n - 1} \\ 0 & \text { otherwise }, \end{array} \right.
\]

以及

\[
f _ {2 n - 1} (x) = \left\{ \begin{array}{l l} \sqrt {2} \sin (2 n \pi x) & \text { if } 0 <   x <   1 / 2 \\ 2 ^ {n} & \text { if } 1 - 2 ^ {- 2 n} <   x <   1 - 2 ^ {- 2 n - 1} \\ 0 & \text { otherwise }, \end{array} \right.
\]

对所有 \(n \geqslant 1\) 和每个 \(x \in X\) 成立。

\begin{enumerate}
\def\labelenumi{\alph{enumi})}
\item
  族 \((f_n)_{n\geqslant 1}\) 在 \(\mathcal{L}^2 (\mathrm{X},\mu)\) 中是规范正交的。
\item
  对每个 \((x,y)\in \mathrm{X}\times \mathrm{X}\)，我们令
\end{enumerate}

\[
\mathrm{N} (x, y) = \sum_ {n = 1} ^ {+ \infty} 2 ^ {- n} (f _ {2 n - 1} (x) f _ {2 n - 1} (y) - f _ {2 n} (x) f _ {2 n} (y)).
\]

证明 \(|\mathrm{N}(x,y)| \leqslant 2^{3/2}\) 对所有 \((x,y) \in \mathrm{X} \times \mathrm{X}\) 成立。由此推出 \(\mathrm{N}_y \in \mathcal{L}^\infty(\mathrm{X})\) 对所有 \(y \in \mathrm{X}\) 成立，且

\[
\int_ {X} ^ {*} \| N _ {y} \| _ {\infty} ^ {2} d y <   + \infty .
\]

\begin{enumerate}
\def\labelenumi{\alph{enumi})}
\setcounter{enumi}{2}
\tightlist
\item
  证明通过令
\end{enumerate}

\[
u _ {\mathrm{N}} (f) (y) = \int_ {\mathrm{X}} \mathrm{N} (x, y) f (x) d \mu
\]

对 \(f \in \mathcal{L}^{1}(X, \mu)\) 和 \(y \in X\) ，通过过渡到商空间，可以定义一个从 \(\mathrm{L}^{1}(X, \mu)\) 到 \(\mathrm{L}^{2}(X, \mu)\) 的连续线性映射。

\begin{enumerate}
\def\labelenumi{\alph{enumi})}
\setcounter{enumi}{3}
\tightlist
\item
  线性映射 \(u_{N}\) 不是紧的。（考虑序列 \(g_{n}=2^{n}(f_{2n-1}-f_{2n})\) ；证明 \(g_{n}\) 可积，且序列 \((u_{\mathrm{N}}(g_{n}))\)（在 \(\mathrm{L}^{2}(\mathrm{X},\mu)\) 中）没有收敛子列。）
\end{enumerate}

\begin{enumerate}
\def\labelenumi{\arabic{enumi})}
\setcounter{enumi}{2}
\tightlist
\item
  设 X 是局部紧拓扑空间，\(\mu\) 是 X 上的一个正测度，\(Y \subset X\) 是局部 \(\mu\) -可忽略的子集。设 N 是 \(X \times Y\) 的特征函数。
\end{enumerate}

\begin{enumerate}
\def\labelenumi{\alph{enumi})}
\tightlist
\item
  我们有
\end{enumerate}

\[
\int_ {X \times X} ^ {*} | N (x, y) f (x) g (y) | d (\mu \otimes \mu) (x, y) = 0
\]

对 \(f, g \in \mathrm{L}^2(\mathrm{X}, \mu)\) 成立。

\begin{enumerate}
\def\labelenumi{\alph{enumi})}
\setcounter{enumi}{1}
\tightlist
\item
  若 \(f \in \mathrm{L}^{2}(\mathrm{X}, \mu) - \mathrm{L}^{1}(\mathrm{X}, \mu)\)，则由 \(y \in X\) 中使 \(x \mapsto \mathrm{N}(x, y)f(x)\) 为 \(\mu\) -可积的那些 y 组成的集合与 Y 重合。
\end{enumerate}

\begin{enumerate}
\def\labelenumi{\arabic{enumi})}
\setcounter{enumi}{3}
\tightlist
\item
  设 \(\alpha \in ]0,1]\) 。设 I 是 \(\mathbf{R}\) 的一个紧区间。我们用 \(\mathcal{C}^{\alpha}(\mathrm{I})\) 表示 \(\mathcal{C}(\mathrm{I})\) 的由满足下列条件的函数 \(f\) 组成的向量子空间：存在 \(\mathrm{C} \geqslant 0\) 满足
\end{enumerate}

\[
| f (x) - f (y) | \leqslant \mathrm{C} | x - y | ^ {\alpha}
\]

对所有 \((x,y)\in\mathrm{I}\times\mathrm{I}\) 成立。

我们赋给 \(\mathcal{C}^{\alpha}(\mathrm{I})\) 范数

\[
\| f\|_{\alpha} = \sup_{x\in \mathrm{I}}|f(x)| + \sup_{\substack{(x,y)\in \mathrm{I}\times \mathrm{I}\\ x\neq y}}\frac{|f(x) - f(y)|}{|x - y|^{\alpha}}.
\]

\begin{enumerate}
\def\labelenumi{\alph{enumi})}
\item
  空间 \(\mathcal{C}^{\alpha}(\mathrm{I})\) 是巴拿赫空间。
\item
  典范单射 \(\mathcal{C}^{\alpha}(\mathrm{I})\to \mathcal{C}(\mathrm{I})\) 是紧的。
\item
  若 \(0 < \alpha < \beta \leqslant 1\)，则典范单射 \(\mathcal{C}^{\beta}(\mathrm{I}) \to \mathcal{C}^{\alpha}(\mathrm{I})\) 是紧的。\\
\item
  索伯列夫空间 \(\mathrm{H}^{1}(\mathring{\mathrm{I}})\)（参见 IV, p. 221, 第 14 目）等同于 \(\mathcal{C}^{1/2}(\mathrm{I})\) 的一个子空间，但典范单射 \(\mathrm{H}^{1}(\mathring{\mathrm{I}}) \to \mathcal{C}^{1/2}(\mathrm{I})\) 不是紧的。
\end{enumerate}

\exercisesection{}

\begin{enumerate}
\def\labelenumi{\arabic{enumi})}
\item
  设 E 和 F 是希尔伯特空间，u 是从 E 到 F 的一个连续线性映射。证明 u 是弗雷德霍姆映射，当且仅当 u 的像是闭的且 u 和 \(u^{*}\) 的核是有限维的。
\item
  设 \(p\) 是满足 \(1 \leqslant p < +\infty\) 的实数。设 \(u\) 是 \(\ell_{\mathrm{K}}^{p}(\mathbf{N})\) 的一个对角自同态。证明 \(u\) 是弗雷德霍姆自同态，当且仅当它是里斯自同态。特别地，\(\ell_{\mathrm{K}}^{p}(\mathbf{N})\) 的每个对角弗雷德霍姆自同态的指标为零。
\item
  \begin{enumerate}
  \def\labelenumii{\alph{enumii})}
  \tightlist
  \item
    存在一个巴拿赫空间 E、E 的一个连续自同构 u 以及 E 的一个里斯自同态 v，使得 \(u \circ v\) 不是里斯自同态。甚至存在这样的例子，其中 \(u \circ v - v \circ u\) 是有限秩的。
  \end{enumerate}
\end{enumerate}

\begin{enumerate}
\def\labelenumi{\alph{enumi})}
\setcounter{enumi}{1}
\tightlist
\item
  存在一个巴拿赫空间 E、E 的一个自同构 u 以及一个有限秩自同态 h，使得 \(u + h\) 不是里斯自同态。
\end{enumerate}

\begin{enumerate}
\def\labelenumi{\arabic{enumi})}
\setcounter{enumi}{3}
\tightlist
\item
  设 E 是无限维复希尔伯特空间。
\end{enumerate}

\begin{enumerate}
\def\labelenumi{\alph{enumi})}
\item
  E 上的酉算子群 \(\mathbf{U}(\mathrm{E})\)，赋以由 \(\mathcal{L}(\mathrm{E})\) 的范数诱导的拓扑，是道路连通的。（证明若 u 是酉的，则存在一个埃尔米特算子 v，使得 \(u = \exp(iv)\) 。）
\item
  E 的线性同构群是道路连通的。（利用极分解，参见 I, p. 139, 定义 3。）
\item
  对每个 \(n \in Z\)，E 的指标为 n 的弗雷德霍姆自同态组成的集合是开且连通的。它是非空的。
\end{enumerate}

\begin{enumerate}
\def\labelenumi{\arabic{enumi})}
\setcounter{enumi}{4}
\tightlist
\item
  设 E 和 F 是无限维巴拿赫空间。连续线性映射 \(u: \mathrm{E} \to \mathrm{F}\) 称为严格奇异的，如果不存在无限维的闭子空间 \(\mathrm{F}_1 \subset \mathrm{E}\)，使得 \(u\) 通过过渡到子空间诱导一个从 \(\mathrm{F}_1\) 到 \(u(\mathrm{F}_1)\) 上的同构。
\end{enumerate}

\begin{enumerate}
\def\labelenumi{\alph{enumi})}
\tightlist
\item
  设 u 是从 E 到 F 的一个连续线性映射。下列条件等价：
\end{enumerate}

\begin{enumerate}
\def\labelenumi{(\roman{enumi})}
\item
  映射 \(u\) 是严格奇异的；
\item
  对 E 的每个无限维闭子空间 \(F_{1}\)，不存在实数 \(\delta > 0\)，使得 \(\|u(x)\| \geqslant \delta\|x\|\) 对所有 \(x \in F_{1}\) 成立。
\end{enumerate}

\begin{enumerate}
\def\labelenumi{\alph{enumi})}
\setcounter{enumi}{1}
\item
  设 \(u \colon \mathrm{E} \to \mathrm{F}\) 是一个连续线性映射，使得 \(u\) 的像是闭的且 \(u\) 的核是有限维的。若 \(v \colon \mathrm{E} \to \mathrm{F}\) 是严格奇异的，则 \(u + v\) 的像是闭的且其核是有限维的。
\item
  E 的严格奇异自同态组成的集合是 \(\mathcal{L}(\mathrm{E})\) 的一个闭双边理想。
\item
  每个紧线性映射都是严格奇异的。若 E 和 F 是希尔伯特空间，则每个严格奇异映射 \(E \rightarrow F\) 都是紧的。
\item
  存在巴拿赫空间 E 和 F 以及一个严格奇异但非紧的映射 \(E \rightarrow F\) 。
\end{enumerate}

\(f)\) 设 \(p\) 和 \(r\) 是满足 \(1 \leqslant p < r\) 的实数。每个连续线性映射 \(\ell_{\mathrm{K}}^{r}(\mathbf{N}) \to \ell_{\mathrm{K}}^{p}(\mathbf{N})\) 都是严格奇异的。（利用 III, p. 106 习题 10。）

\begin{enumerate}
\def\labelenumi{\arabic{enumi})}
\setcounter{enumi}{5}
\item
  设 E 是复希尔伯特空间。设 u 是 E 的一个弗雷德霍姆自同态。若 u 是正规的，则 \(\text{ind}(u) = 0\) 。
\item
  设 E 和 F 是巴拿赫空间，u 是 \(\mathcal{L}(\mathrm{E};\mathrm{F})\) 的一个元素。证明 u 是弗雷德霍姆映射，当且仅当下列条件被满足：
\end{enumerate}

\begin{enumerate}
\def\labelenumi{(\roman{enumi})}
\tightlist
\item
  存在 E 上的一个连续半范数 p 和 \(C \geqslant 0\)，使得从 E 到 \((E, p)\) 的分离完备化的典范映射是紧的，且
\end{enumerate}

\[
\| x \| \leqslant \mathrm{C} \| u (x) \| + p (x)
\]

对所有 \(x \in E\) 成立；

\begin{enumerate}
\def\labelenumi{(\roman{enumi})}
\setcounter{enumi}{1}
\tightlist
\item
  存在一个连续半范数 \(q\)（\(\mathrm{F}'\) 上的）和 \(\mathrm{C}' \geqslant 0\)，使得从 \(\mathrm{F}'\) 到 \((\mathrm{F}', q)\) 的分离完备化的典范映射是紧的，且
\end{enumerate}

\[
\| y ^ {\prime} \| \leqslant \mathrm{C} ^ {\prime} \| ^ {t} u (y ^ {\prime}) \| + q (y ^ {\prime})
\]

对所有 \(y' \in F'\) 成立。

\begin{enumerate}
\def\labelenumi{\arabic{enumi})}
\setcounter{enumi}{7}
\tightlist
\item
  设 U 是 C 的一个连通开子集。我们用 \(\mathcal{H}(\mathrm{U})\) 表示 U 上的赋有紧收敛拓扑的全纯函数空间（EVT III, p. 10, 例 \(d\)）。
\end{enumerate}

\begin{enumerate}
\def\labelenumi{\alph{enumi})}
\item
  设 \(a \in \mathcal{H}(U)\) 。乘以 a 是 \(\mathcal{H}(U)\) 的一个弗雷德霍姆自同态，当且仅当 a 的零点集是有限的。然后计算它的指标。
\item
  设 \(\mathcal{U} = (\mathrm{U}_{i})_{i \in \mathrm{I}}\) 是 U 的一个由非空连通开集组成的覆盖。我们令
\end{enumerate}

\[
\begin{array}{l} \mathrm{I} _ {2} = \{(i, j) \in \mathrm{I} \times \mathrm{I} | \mathrm{U} _ {i} \cap \mathrm{U} _ {j} \neq \varnothing \}, \\ \mathrm{I} _ {3} = \{(i, j, k) \in \mathrm{I} \times \mathrm{I} \times \mathrm{I} | \mathrm{U} _ {i} \cap \mathrm{U} _ {j} \cap \mathrm{U} _ {k} \neq \varnothing \}. \end{array}
\]

我们考虑向量空间

\[
\begin{array}{l} \mathrm{Z} ^ {1} (\mathcal {U}; \mathbf {C}) = \{(c _ {i, j}) \in \mathbf {C} ^ {\mathrm{I} _ {2}} | c _ {i j} + c _ {j k} + c _ {k i} = 0 \text { for all } (i, j, k) \in \mathrm{I} _ {3} \} \\ \mathrm{B} ^ {1} (\mathcal {U}; \mathbf {C}) = \{(c _ {i, j}) \in \mathbf {C} ^ {\mathrm{I} _ {2}} | \text { there exists } c ^ {\prime} \in \mathbf {C} ^ {\mathrm{I}} \text { such that } \end{array}
\]

\[
c _ {i j} = c _ {i} ^ {\prime} - c _ {j} ^ {\prime} \text {   for all   } (i, j) \in \mathrm{I} _ {2} \}.
\]

我们有 \(\mathrm{B}^1 (\mathcal{U};\mathbf{C})\subset \mathrm{Z}^1 (\mathcal{U};\mathbf{C})\) 。我们令

\[
\mathrm{H} ^ {1} (\mathcal {U}; \mathbf {C}) = \mathrm{Z} ^ {1} (\mathcal {U}; \mathbf {C}) / \mathrm{B} ^ {1} (\mathcal {U}; \mathbf {C}).
\]

设 \(\mathcal{V}=(\mathrm{V}_{\alpha})_{\alpha\in\mathrm{A}}\) 是 U 的一个由非空连通开集组成且比 U 更细的覆盖，并设 \(\varrho\colon A\to I\) 是一个映射，使得 \(V_{\alpha}\subset U_{\varrho(\alpha)}\) 对所有 \(\alpha\in A\) 成立。

映射 \((c_{ij})_{i,j} \mapsto (c_{\varrho(\alpha)\varrho(\beta)})_{\alpha,\beta}\) 诱导一个不依赖于 \(\varrho\) 的选取的线性映射，从 \(\mathrm{H}^1(\mathcal{U};\mathbf{C})\) 到 \(\mathrm{H}^1(\mathcal{V};\mathbf{C})\) 。

我们用 \(H^{1}(U; \mathbf{C})\) 表示当 U 取遍 U 的由非空连通开集组成的覆盖时所得到的赋有这些线性映射的向量空间 \(H^{1}(\mathcal{U}; \mathbf{C})\) 的归纳极限。

\begin{enumerate}
\def\labelenumi{\alph{enumi})}
\setcounter{enumi}{2}
\tightlist
\item
  设 D 是 \(\mathcal{H}(\mathrm{U})\) 的由 \(Df = f'\) 定义的自同态。证明存在一个从 D 的余核到 \(\mathrm{H}^1 (\mathrm{U};\mathbf{C})\) 上的向量空间同构。由此推出 D 是弗雷德霍姆自同态，当且仅当空间 \(\mathrm{H}^1 (\mathrm{U};\mathbf{C})\) 是有限维的。然后计算 D 的指标。
\end{enumerate}

\exercisesection{}

\begin{enumerate}
\def\labelenumi{\arabic{enumi})}
\item
  设 E 和 F 是巴拿赫空间，u 是一个连续线性映射 \(E \rightarrow F\) 。我们有 \(((u)) > 0\)，当且仅当 u 的像在 F 中是闭的。
\item
  设 \(p \in [1, +\infty]\) ，且 \(\mathrm{E} = \ell_{\mathbf{C}}^{p}(\mathbf{N})\) 。设 u 是 E 的一个对角自同态。对每个 \(\lambda \in C\) ，用 u 的谱计算 \(((u - \lambda 1_{\mathrm{E}}))\) 。
\item
  设 \(p \in [1, +\infty]\) 且 \(\mathrm{E} = \ell_{\mathrm{K}}^{p}(\mathbf{N})\)，并设 \(u\) 是移位自同态
\end{enumerate}

\[
(x _ {0}, x _ {1}, \dots ,) \mapsto (0, x _ {0}, x _ {1}, \dots)
\]

（E 上的）。证明 u 是指标为 1 的弗雷德霍姆自同态，且 \(((u)) = 1\) 。

\begin{enumerate}
\def\labelenumi{\arabic{enumi})}
\setcounter{enumi}{3}
\item
  设 E 是巴拿赫空间。映射 \(u \mapsto ((u))\)（从 \(\mathcal{L}(\mathrm{E})\) 到 R）是连续的，当且仅当 E 的维数 \(\leqslant 1\) 。
\item
  设 E 是希尔伯特空间，\(u \in \mathcal{L}(\mathrm{E})\) 。我们有 \(((u)) = \inf(\mathrm{Sp}(|u|) - \{0\})\)（\(|u|\) 的定义参见 I, p. 139, 定义 3）。
\item
  设 \(E = \ell^{2}(\mathbf{N})\)，\((e_{n})_{n \in \mathbf{N}}\) 是 E 的规范正交基。设 u 是 E 的对角自同态，满足 \(u(e_{0}) = e_{0}\) 且 \(u(e_{n}) = 2e_{n}\)（若 \(n \geqslant 1\)）。证明 \(((u)) = 1\)，且 \(u + h\) 对每个满足 \(\|h\| < 2\) 的 h 是弗雷德霍姆自同态。
\item
  设 E 和 F 是巴拿赫空间，u 是 \(\mathcal{Q}\mathcal{M}(\mathrm{E};\mathrm{F})\)（相应地，\(\mathcal{Q}\mathcal{E}(\mathrm{E};\mathrm{F})\)）的一个元素。设 h 是 \(\mathcal{L}(\mathrm{E};\mathrm{F})\) 的一个元素。假设存在正实数 a 和 b，使得 \(a + b((u))^{-1} < 1\)，从而
\end{enumerate}

\[
\| h (x) \| \leqslant a \| u (x) \| + b \| x \|\tag{2}
\]

对所有 \(x \in E\) 成立。

证明 \(u + h \in \mathcal{Q}\mathcal{M}(\mathrm{E};\mathrm{F})\)（相应地，\(u + h \in \mathcal{Q}\mathcal{E}(\mathrm{E};\mathrm{F})\)）且

\[
\begin{array}{c} \dim \operatorname{Ker} (u + h) \leqslant \dim \operatorname{Ker} (u), \\ \dim \operatorname{Coker} (u + h) \leqslant \dim \operatorname{Coker} (u), \\ \operatorname{ind} (u + h) = \operatorname{ind} (u). \end{array}
\]

（化归到 \(a\) 和 \(b\) 都 \(>0\) 的情形，并设 \(c > 0\) 满足 \(a + b/c < 1\) ；考虑 E 上由下式定义的范数 \(p\)

\[
p (x) = a \| u (x) \| + b \| x \| \quad \text {   for   } x \in \mathrm{E},
\]

并设 \(E_{1}\) 是赋有范数 p 的巴拿赫空间 E；我们有 \(\mathcal{Q}\mathcal{M}(E_{1};F)=\mathcal{Q}\mathcal{M}(E;F)\)，\(\mathcal{Q}\mathcal{E}(E_{1};F)=\mathcal{Q}\mathcal{E}(E;F)\) 且 \(\mathcal{L}(E_{1};F)=\mathcal{L}(E;F)\) 。然后应用 III, p. 62 的引理 2。）

\begin{enumerate}
\def\labelenumi{\arabic{enumi})}
\setcounter{enumi}{7}
\item
  设 E 是希尔伯特空间。不用博苏克–乌拉姆定理，直接证明 III, p. 64 的定理 4。
\item
  设 \(r \geqslant 1\) 是整数。若 \(f: \mathrm{M}_1 \to \mathrm{M}_2\) 是 \(\mathrm{C}^r\) 类流形之间的一个 \(\mathrm{C}^r\) 类态射，则称 \(x \in \mathrm{M}_1\) 是 \(f\) 的临界点，如果 \(f\) 在 \(x\) 处不是淹没（VAR, R1, 5.9.1）。设 X 是 \(f\) 的临界点集。集合 \(f(\mathrm{X})\)（在 \(\mathrm{M}_2\) 中）称为 \(f\) 的临界值集，而 \(f(\mathrm{M}_1) = f(\mathrm{X})\) 称为 \(f\) 的正则值集。
\end{enumerate}

若 \(M_{1}\) 和 \(M_{2}\) 是 \(C^{r}\) 类实流形（VAR, R1, 5.1.5, p. 35），流形的一个态射 \(f: M_{1} \to M_{2}\)（\(C^{r}\) 类）称为弗雷德霍姆的，如果对每个 \(x \in M_{1}\) ，切线性映射 \(\mathrm{T}_{x}(f)\) 是从巴拿赫空间 \(\mathrm{T}_{x}(\mathrm{M}_{1})\) 到巴拿赫空间 \(\mathrm{T}_{f(x)}(\mathrm{M}_{2})\) 的弗雷德霍姆自同态。a) 设 \(f: M_{1} \to M_{2}\) 是一个弗雷德霍姆态射。若 \(M_{1}\) 是连通的，则指标 \(\operatorname{ind}(\mathrm{T}_{x}(f))\) 不依赖于 \(x \in M_{1}\) 。它称为弗雷德霍姆态射 f 的指标，记作 \(\operatorname{ind}(f)\) 。

\begin{enumerate}
\def\labelenumi{\alph{enumi})}
\setcounter{enumi}{1}
\tightlist
\item
  设 U 是巴拿赫空间 E 的一个开子集，满足 \(0 \in U\)，并设 F 是巴拿赫空间。设 \(f: U \to F\) 是一个 \(C^{r}\) 类弗雷德霍姆态射，并设 \(u = \mathrm{T}_{0}(f)\) 。存在一个巴拿赫空间 \(E_{1}\)、U 的一个卡 \((\mathrm{U}_{1}, \varphi, \mathrm{E}_{1} \times \mathrm{Ker}(u))\)（满足 \(0 \in U_{1}\)）、F 的一个卡 \((\mathrm{U}_{2}, \psi, \mathrm{E}_{1} \times \mathrm{Coker}(u))\)（满足 \(f(0) \in U_{2}\)），以及一个弗雷德霍姆态射 \(g: \varphi(\mathrm{U}_{1}) \to \mathrm{E}_{1} \times \mathrm{Coker}(u)\)（\(C^{r}\) 类），使得
\end{enumerate}

\[
(\psi \circ f \circ \varphi^ {- 1}) (a, b) = (a, g (a, b))
\]

对所有 \((a,b)\in \varphi (\mathrm{U}_1)\subset \mathrm{E}_1\times \mathrm{Ker}(u).\) 成立。

\begin{enumerate}
\def\labelenumi{\alph{enumi})}
\setcounter{enumi}{2}
\item
  设 \(f \colon \mathrm{M}_1 \to \mathrm{M}_2\) 是 \(\mathbf{C}^r\) 类流形之间的一个 \(\mathbf{C}^r\) 类弗雷德霍姆态射。对每个 \(y \in \mathbf{M}_2\)，存在一个开邻域 \(\mathrm{V}\)（\(y\) 的），使得 \(f\) 通过过渡到子空间定义一个逆紧映射 \(f^{-1}(\mathrm{V}) \to \mathrm{V}\) 。
\item
  设 U 是巴拿赫空间 E 的一个包含 0 的开子集，并设 F 是巴拿赫空间。设 \(f \colon \mathrm{U} \to \mathrm{F}\) 是一个 \(\mathrm{C}^{\infty}\) 类弗雷德霍姆态射。对每个 \(x \in \mathrm{U}\) 和 F 中 \(f(x)\) 的每个邻域 V，存在 \(f\) 在 V 中的一个正则值。（利用第二萨德定理，VAR, R2, 10.1.3, p. 37。）e) 设 \(f \colon \mathrm{M}_1 \to \mathrm{M}_2\) 是 \(\mathrm{C}^{\infty}\) 类流形的一个 \(\mathrm{C}^{\infty}\) 类弗雷德霍姆态射。若 \(\mathrm{M}_1\) 和 \(\mathrm{M}_2\) 具有可数基，则 \(f\) 的临界值集是贫集（“斯梅尔定理”）。
\end{enumerate}

\(f)\) 设 \(f\colon \mathrm{M}_1\to \mathrm{M}_2\) 是具有可数基的 \(\mathrm{C}^{\infty}\) 类流形的一个 \(\mathrm{C}^{\infty}\) 类弗雷德霍姆态射。若 \(\operatorname {ind}(f) < 0\)，则 \(f\) 的像的内部是空的。

\(g)\) 设 \(f\colon \mathrm{M}_1\to \mathrm{M}_2\) 是具有可数基的 \(\mathrm{C}^{\infty}\) 类流形的一个 \(\mathrm{C}^{\infty}\) 类弗雷德霍姆态射。对每个位于一个贫集之外的 \(y\in \mathrm{M}_2\)，集合 \(f^{-1}(y)\) 是维数为 \(\operatorname {ind}(f)\) 的纯子流形。